\documentclass[10pt,journal,compsoc]{IEEEtran}
\usepackage[utf8]{inputenc}
\usepackage[T1]{fontenc}
\usepackage{url}
\usepackage[hidelinks]{hyperref}
\usepackage{cite}
\usepackage{amsmath,amssymb,amsfonts}
\usepackage{algorithmic}
\usepackage{graphicx}
\usepackage{textcomp}
\usepackage{xcolor}
\usepackage{tabularx}
\usepackage{listings}
\usepackage{array}

\lstset{
  basicstyle=\small\ttfamily,
  breaklines=true,
  frame=single,
  columns=fullflexible
}

\begin{document}

\title{Paper}
\author{Bryan Alexander Buchorn \\ Independent Researcher \\ bryan.alexander@buchorn.com}
\markboth{CEREBRUM v1.2.4 Hardened Edition · March 2026}{}

\maketitle

\begin{abstract}
This module forms a critical component of the CEREBRUM framework, a community-structured graph attention engine for Knowledge Graph reasoning. It focuses on Framework Implementation Guide.
\end{abstract}

\begin{IEEEkeywords}
Knowledge Graph, Graph Attention, DSCF, CSA, Hierarchical Reasoning.
\end{IEEEkeywords}

\section{CEREBRUM: Community-Structured Graph Attention for Knowledge Graph Reasoning}

\textbf{Authors}: Bryan Alexander Buchorn
\textbf{Affili\-ations}: Independent Researcher, Las Vegas, NV, USA
\textbf{Contact}: bryan.alexander@buchorn.com
\textbf{Date}: March 2026
\textbf{Status}: Version 1.2.0 $\cdot$ Phase 21 COMPLETE — 130 advanced tests passing, all E2E journeys passing
\textbf{License}: Proprietary — all rights reserved

---

\subsection{Quick Reference}

\textbf{What it is}: A framework that lets a Knowledge Graph reason like a Transformer —
using community structure as attention heads, BFS hop depth as layer depth, and
graph-structural features as positional encodings. No training required. No LLM required.
Every inference step is a verifiable graph edge.

\textbf{The two core algorithms}:

\begin{equation}
\text{DSCF (community detection):}
\end{equation}
For each node $v$ at each iteration:
\begin{itemize}
\item $\text{lpa\_cid} = \text{majority vote among neighbors (local signal)}$
\item $\text{mod\_cid} = \text{best modularity gain } \Delta Q \text{ (global signal)}$
\item If signals agree and $\neq \text{current}$: \textbf{MOVE (anchor)}
\item Elif disagree: \textbf{MOVE} with probability weighted by $(\text{lpa\_conf} \times \tau)$ vs $(\text{mod\_conf} \times (2-\tau))$

\end{itemize}
Temperature schedule: $\tau_{t+1} = \max(\tau_t \times 0.92, 0.01)$

\begin{equation}
\text{CSA (attention weight for edge } u \to v \text{ at hop } k):
\end{equation}
\begin{equation}
a(u,v,k) = \sigma( \alpha \cdot \text{sim}(u,v) + \beta \cdot S_{com}(u,v) + \gamma \cdot w_{rel} - \delta \cdot d_{norm} + \epsilon \cdot \phi(k) )
\end{equation}
Defaults: $\alpha=0.4, \beta=0.4, \gamma=0.1, \delta=0.05, \epsilon=0.05$
$w_{rel}$: Metaedge Bridge Bonus (default 0.0, recommended 0.4 for inter-type reasoning)

\textbf{Transformer $\to$ KG mapping}:

\begin{center}
\begin{tabularx}{\columnwidth}{|>{\raggedright\arraybackslash}X|>{\raggedright\arraybackslash}X|}
\hline
Transformer & CEREBRUM \\ \hline
Attention head & DSCF community \\ \hline
Layer depth & BFS hop count \\ \hline
Positional encoding & PageRank + betweenness + degree \\ \hline
Attention weight & CSA formula above \\ \hline
Context window & Ego-network radius R \\ \hline
KV cache & Materi\-alized path store \\ \hline
\end{tabularx}
\end{center}

\textbf{System component names:}

\begin{center}
\begin{tabularx}{\columnwidth}{|>{\raggedright\arraybackslash}X|>{\raggedright\arraybackslash}X|}
\hline
Name & Role \\ \hline
\textbf{CEREBRUM} & The overarching product/framework \\ \hline
\textbf{THALAMUS} & Ingestion engine — adapters, embedding, structural encoding, STDP discretizer \\ \hline
\textbf{CORTEX} & Core reasoning engine — DSCF + CSA + BeamTraversal + AnswerExtractor \\ \hline
\textbf{REM Engine} & Graph self-reorg\-anization — prune/consolidate/synthesize \\ \hline
\textbf{Bridge Twin Engine} & Experience-dependent structural relay nodes \\ \hline
\end{tabularx}
\end{center}

\textbf{Repo layout} (target structure for standalone project):

\begin{lstlisting}
parallax/
├── adapters/      networkx, neo4j, rdf, csv, file_adapter, stream_adapter    [THALAMUS]
├── core/
│   ├── embedding_engine, structural_encoder, discretizer                      [THALAMUS]
│   ├── community_engine, attention_engine, parameter_learner, kge_engine      [CORTEX]
│   ├── rem_engine, bridge_engine                                               [REM / Bridge Twin]
│   ├── insight_engine, insight_validator, meta_insight_engine                 [Verification / Metacognition]
│   └── graph_adapter, hardware, security, contradiction_engine
├── reasoning/     traversal, path_scorer, answer_extractor                    [CORTEX]
├── llm_bridge/    context_formatter
├── api/           server, schemas (+ /stream/* endpoints)
├── cli/           parallax.py
├── tests/         test_dscf, test_csa, test_traversal, fixtures/toy_graph.csv
├── benchmarks/    webqsp_eval, metaqa_eval, baseline_comparison
├── examples/      wikidata_quickstart, neo4j_quickstart, csv_quickstart, Validation_Walkthrough.ipynb
├── pyproject.toml
├── README.md
└── PAPER.md       (this file)
\end{lstlisting}

\textbf{Current phase}: Phase 21 complete (v1.2.0). All core E2E journeys and 130 advanced tests passing. All Phase 18-20 features are stable, and the SignalEncoder alignment has been fixed.

---

\subsection{Abstract}

We propose \textbf{CEREBRUM}, a novel framework that enables Knowledge Graphs (KGs)
to perform multi-hop reasoning using the same structural principles that make
Transformer-based Large Language Models powerful — without requiring an LLM,
without training data, and with full inter\-pretability of every inference step.

The central contr\-ibution is \textbf{Community-Structured Attention (CSA)}: a
mechanism in which graph communities serve as attention heads, graph traversal
replaces matrix multi\-plication, and hop depth replaces layer depth. Unlike
Graph Attention Networks (GATs), which apply learned attention within hard
adjacency constraints, CSA uses community membership as a soft global
constraint that captures both local topological cohesion and global structural
signi\-ficance simul\-taneously.

\subsubsection{Table 1: Comparative Value Proposition}

\begin{center}
{\footnotesize
\begin{tabularx}{\columnwidth}{|>{\raggedright\arraybackslash}X|>{\raggedright\arraybackslash}X|>{\raggedright\arraybackslash}X|>{\raggedright\arraybackslash}X|}
\hline
Feature & Standard RAG & GraphRAG (Microsoft) & CEREBRUM \\ \hline
\textbf{Primary Reasoner} & LLM & LLM & \textbf{Knowledge Graph} \\ \hline
\textbf{Logic Source} & Probab\-ilistic weights & LLM-generated summaries & \textbf{Graph Topology (DSCF/CSA)} \\ \hline
\textbf{Halluc\-ination Risk} & High & Medium & \textbf{Zero (Grounded Paths)} \\ \hline
\textbf{Interp\-retability} & None (Black-box) & Medium (Text chunks) & \textbf{Absolute (Verifiable Edges)} \\ \hline
\textbf{Context Window} & Limited by Token Count & Limited by Chunk Count & \textbf{Scale-Invariant (Beam Search)} \\ \hline
\end{tabularx}
}
\end{center}

This is made possible by a second contr\-ibution: the \textit{}Dual-Signal Community
Fusion (DSCF)\textit{} algorithm, which produces communities that encode both LPA
majority-vote structure (local) and modularity gain (global) in a single
partition. We show that DSCF communities possess a structural duality that
maps naturally to the dual character of multi-head attention in Transf\-ormers.

Together, CSA and DSCF form an archi\-tecture where a KG can answer multi-hop
questions by traversing itself, with every reasoning step grounded in explicit
graph edges, every conclusion traceable to a path, and no LLM required for
inference — though one may optionally be used for natural language generation.

---

\subsection{1. Introd\-uction}

\subsubsection{1.1 The Gap Between Knowledge Graphs and Language Models}

Knowledge Graphs and Large Language Models represent two funda\-mentally
different approaches to knowledge repre\-sentation and reasoning.

Knowledge Graphs store knowledge explicitly: entities as nodes, relat\-ionships
as typed edges, facts as (subject, predicate, object) triples. This makes
them precise, verifiable, and updatable without retraining. However, they
cannot reason beyond what is explicitly stored. Multi-hop inference — "Marie
Curie discovered Polonium, Polonium is radioactive, therefore Marie Curie
discovered a radioactive element" — requires either hardcoded graph traversal
queries or external reasoning systems.

Large Language Models store knowledge implicitly in billions of weight
parameters. They can reason, generalize, and synthesize across domains. But
this implicit repre\-sentation is opaque, cannot be updated without expensive
fine-tuning, and is prone to hallu\-cination — generating plausible-sounding
but incorrect facts with no mechanism for ground-truth verif\-ication.

The field has responded with hybrid approaches: Retrieval-Augmented Generation
(RAG), Knowledge-Graph-augmented LLMs, and GraphRAG. In all of these, the KG
is a retrieval store and the LLM does the reasoning. The KG remains passive.

\textbf{CEREBRUM inverts this relat\-ionship.} The KG reasons. The LLM, if present,
only generates natural language from the KG's output. Every inference step is
a graph traversal, every conclusion is a path, and every path can be verified.

\subsubsection{1.2 The Core Observation}

A Transformer's power comes from three mechanisms working together:

\begin{enumerate}
\item \textbf{Multi-head attention}: different heads specialize on different
\end{enumerate}
   relational aspects of the input (syntactic, semantic, long-range, etc.)
\begin{enumerate}
\item \textbf{Deep composition}: each layer builds on the previous, allowing complex
\end{enumerate}
   multi-step reasoning
\begin{enumerate}
\item \textbf{Positional awareness}: the model knows where each token sits relative
\end{enumerate}
   to others

We observe that Knowledge Graphs have natural analogs for all three:

\begin{enumerate}
\item \textbf{Community structure} serves the role of attention heads: nodes within
\end{enumerate}
   a community have strong mutual relevance; communities specialize on
   different conceptual domains.
\begin{enumerate}
\item \textbf{BFS hop depth} serves the role of layer depth: each hop is one step of
\end{enumerate}
   composed reasoning.
\begin{enumerate}
\item \textbf{Graph-structural features} (PageRank, betweenness, degree) serve the
\end{enumerate}
   role of positional encoding: they tell the model where each entity sits in
   the global information landscape.

The question is: can these analogs be made operational — not merely
metap\-horical? We argue yes, and demonstrate the archi\-tecture to do so.

\subsubsection{1.3 Contri\-butions}

This paper makes three primary contr\-ibutions:

\begin{enumerate}
\item \textbf{The CEREBRUM archi\-tecture}: a complete mapping of Transformer components
\end{enumerate}
   to KG operations, enabling multi-hop reasoning via graph traversal alone.

\begin{enumerate}
\item \textbf{Community-Structured Attention (CSA)}: a novel attention mechanism using
\end{enumerate}
   community membership as a soft global constraint on graph traversal,
   bridging the gap between local GAT-style attention and global Transformer-
   style attention.

\begin{enumerate}
\item \textbf{Dual-Signal Community Fusion (DSCF)}: a novel community detection
\end{enumerate}
   algorithm combining LPA majority-vote and modularity gain simul\-taneously
   at each node update, producing communities with dual short-range/long-range
   character that maps to multi-head attention's dual speci\-alization.

---

\subsection{2. Background and Related Work}

\subsubsection{2.1 Graph Neural Networks}

Graph Neural Networks (GNNs) [Scarselli et al., 2009] generalize neural
networks to graph-structured data. The message-passing paradigm [Gilmer et
al., 2017] defines node updates as aggre\-gations of neighbor repre\-sentations.

Graph Attention Networks (GATs) [Velickovic et al., 2018] introduce learned
attention weights between connected nodes. Key limitations for our purposes:
(a) attention is restricted to direct neighbors — no global context;
(b) communities are not considered; (c) training labels required.

GraphSAGE [Hamilton et al., 2017] samples and aggregates local neigh\-borhoods
but similarly lacks global structural awareness.

None of these use community structure as an organ\-izational principle for
attention.

\subsubsection{2.2 Knowledge Graph Reasoning}

Early KG reasoning systems used rule-based approaches (AMIE [Galarraga et al.,
2013]) or embedding-based methods (TransE [Bordes et al., 2013], RotatE [Sun
et al., 2019]). These methods produce entity embeddings but do not perform
multi-hop traversal in the attention-mechanism sense.

Path-based reasoning approaches (DeepPath [Xiong et al., 2017], MINERVA [Das
et al., 2018]) use reinf\-orcement learning to find paths. These are closer to
our goal but require training data and do not use community structure.

\subsubsection{2.3 LLM + KG Hybrid Systems}

KG-GPT [Yao et al., 2023], KGPT [Chen et al., 2020], and related systems
connect KGs to LLMs as retrieval stores. The LLM always performs the
reasoning step; the KG is passive.

GraphRAG [Edge et al., 2024] uses community detection (Leiden) to summarize
graph clusters as text, which is then passed to an LLM for RAG. This is the
closest existing work. However:
\begin{itemize}
\item Communities are summarized as static text chunks, not used as attention heads
\item The LLM performs all reasoning
\item Paths are not returned or made inter\-pretable
\item The system is not graph-agnostic (designed for Microsoft's pipeline)

\end{itemize}
RAPTOR [Sarthi et al., 2024] builds hiera\-rchical text clusters for RAG but
operates on documents, not graphs, and uses tree structure rather than
community structure.

\subsubsection{2.4 Community Detection}

The Louvain algorithm [Blondel et al., 2008] optimizes modularity greedily
but can produce internally disco\-nnected communities. Leiden [Traag et al.,
2019] fixes this with a refinement phase promoting internal conne\-ctivity.
Label Propagation [Raghavan et al., 2007] is fast and unsup\-ervised but
non-deter\-ministic and produces variable quality.

\textbf{DSCF} (this work) combines LPA and Leiden signals simul\-taneously at each
node update, using a temperature-annealing schedule. See Section 4.2.

\textbf{DSCF is non-deter\-ministic} (inherits LPA's shuffle-order sensitivity).
For repro\-ducible results: run 3-5 trials and select the partition with
highest modularity score. For production: seed the RNG and document the seed.

---

\subsection{3. The Structural Equivalence}

We establish a complete operational mapping between Transformer components
and KG operations. This is not analogy — each mapping is functional.

\begin{center}
{\footnotesize
\begin{tabularx}{\columnwidth}{|>{\raggedright\arraybackslash}X|>{\raggedright\arraybackslash}X|>{\raggedright\arraybackslash}X|}
\hline
Transformer Component & CEREBRUM (KG) Equivalent & Notes \\ \hline
Token & Entity or Relation & Atomic unit of information \\ \hline
Vocabulary & Entity type taxonomy & Closed set of possible types \\ \hline
Token embedding & Entity embedding (TransE/RotatE) & Dense vector per entity \\ \hline
Positional encoding & PageRank + betweenness + degree & Where entity sits globally \\ \hline
Attention head & Community cluster (TSC) & Specialized relational context \\ \hline
Attention weight & CSA weight formula & Sim + community + edge + distance \\ \hline
Context window & Ego-network radius R & How far to traverse \\ \hline
Layer depth L & BFS hop count & Reasoning step count \\ \hline
Wormhole Attention & Federated Link & Cross-graph attention jump \\ \hline
Holographic Index & Compressed Signature & Privacy-preserving discovery \\ \hline
Feed-forward sublayer & Entity-type projection & Type-specific trans\-formation \\ \hline
Residual connection & Previous-hop embedding & Prevents information loss \\ \hline
Layer norma\-lization & Embedding norma\-lization & Prevents value explosion \\ \hline
Output projection & Path decoder / ranker & Maps traversal to answer \\ \hline
KV cache & Materi\-alized path store & Reuse traversal across queries \\ \hline
\end{tabularx}
}
\end{center}

This equivalence has a critical implication: \textit{}the number of attention heads is
not a hyper\-parameter in CEREBRUM — it is determined by the graph's own community
structure.\textit{} A graph with 12 natural communities has 12 attention heads. A graph
with 200 communities has 200. The archi\-tecture adapts to the data.

---

\subsection{4. The CEREBRUM Archit\-ecture}

\subsubsection{4.1 Community-Structured Attention (CSA)}

CSA computes attention weights for graph traversal that incorporate both local graph topology and global community structure.

\textbf{Attention weight formula:}

For entity $u$ attending to entity $v$ at traversal hop $k$:

\begin{equation}
a(u, v, k) = \sigma\left( \alpha \cdot \cos(\vec{e}_u, \vec{e}_v) + \beta \cdot S_{com}(u, v) + \gamma \cdot w_{rel} - \delta \cdot d_{norm}(u, v) + \epsilon \cdot \phi(k) \right)
\end{equation}

Where:
\begin{itemize}
\item $\vec{e}$ is the entity embedding.
\item $S_{com}(u, v)$:
\item $1.0$ if $\text{community}(u) == \text{community}(v)$
\item $0.5$ if communities are adjacent
\item $\exp(-\lambda \cdot d_{com})$ otherwise
\item $w_{rel}$: weight per relation type.
\item $d_{norm}$: normalized shortest path length.
\item $\phi(k)$: hop decay (e.g., $1 / (1 + k)$).
\item $\sigma$: sigmoid activation.


\end{itemize}
\textbf{Default parameter values (zero-shot deployment):}
\begin{itemize}
\item $\alpha$ = 0.4 (embedding similarity)
\item $\beta$ = 0.4 (community membership)
\item $\gamma$ = 0.1 (edge type)
\item $\delta$ = 0.05 (distance penalty)
\item $\epsilon$ = 0.05 (hop decay)

\end{itemize}
Parameters can be learned from (query, answer) pairs for supervised settings.

\textbf{community\_score definition (complete):}

\begin{lstlisting}
community_score(u, v):
  if community(u) == community(v):          return 1.0
  if communities_are_adjacent(u, v):        return 0.5
  else:
    d = community_distance(u, v)
    return exp(-λ · d)

community_distance(u, v):
  # Shortest path (in hops) between the community of u and community of v
  # in the community-level graph, where communities are nodes and two
  # communities are adjacent if ≥1 cross-community edge exists between them.
  # Precomputed once after DSCF converges via BFS on the community graph.
  # λ = 0.5 default (controls cross-community decay rate)
\end{lstlisting}

\textbf{Why this is not a GAT:}

GATs compute \texttt{a(u, v) = f(Wu · emb(u), Wv · emb(v))} — purely from learned
weights on adjacent node pairs. They cannot express the community membership
term $\beta$ $\cdot$ community\_score(u, v), which introduces global structural awareness
without requiring the full O(n²) attention of Transf\-ormers.

CSA is O(n $\cdot$ k̄ $\cdot$ C) where k̄ is average degree and C is the average number
of community-adjacent entities to consider — far cheaper than Transformer
attention while capturing global structure via community membership.

\subsubsection{4.2 Dual-Signal Community Fusion (DSCF)}

DSCF is the community detection algorithm that produces the attention head
structure. It is a key contr\-ibution in its own right.

\textbf{The core innovation:} at each individual node update, both the LPA
majority-vote signal (local topology) and the modularity gain signal (global
structure) are computed. The decision incor\-porates both simul\-taneously,
governed by a temperature parameter:

\begin{lstlisting}
For each node v at each iteration:

  1. LPA signal:
     lpa_cid = argmax over neighbor labels (majority vote)
     lpa_conf = vote_count[lpa_cid] / total_neighbors  ∈ [0, 1]

  2. Modularity signal:
     For each candidate community $C$ adjacent to $v$:
     $$\Delta Q(v \to C) = \frac{k_{v,C}}{m} - \gamma \frac{k_v \sum k_C}{2m^2}$$
     $\text{best\_mod\_cid} = \text{argmax } \Delta Q$
     $\text{mod\_conf} = \min(\text{best\_}\Delta Q \cdot m, 1.0)$

  3. Decision:
     if lpa_cid == best_mod_cid ≠ current:
       MOVE (consensus anchor — high confidence)
     elif both say STAY:
       STAY
     elif only LPA says MOVE:
       MOVE with probability lpa_conf × temperature
     elif only modularity says MOVE:
       MOVE with probability mod_conf × (1 + (1 − temperature))
     else (disagree on different targets):
       lpa_weight = lpa_conf × temperature
       mod_weight = mod_conf × (2 − temperature)
       MOVE to weighted-random choice

  temperature schedule: τ_{t+1} = max(τ_t × cooling, 0.01)
\end{lstlisting}

Post-convergence: Leiden-style conne\-ctivity check — split any community whose
induced subgraph is disco\-nnected.

\textbf{Why DSCF communities are the right attention heads:}

In a trained Transformer:
\begin{itemize}
\item Some heads specialize on local structure (syntactic patterns, adjacent tokens)
\item Other heads specialize on long-range structure (coreference, semantic themes)
\item The combination allows both local and global reasoning simul\-taneously

\end{itemize}
DSCF communities exhibit exactly this dual character. The LPA component supports
communities are locally coherent (nodes in the same community are topol\-ogically
close). The modularity component supports communities are globally significant
(they represent struc\-turally distinct regions of the graph). A node that is in
a DSCF community is there because \textit{both} local and global signals agreed.

This dual property has not previously been used as a basis for attention
mechanisms in any published graph learning system.

\textbf{Comparison to Leiden-only communities:}

Leiden optimizes purely for modularity. On sparse regions of a graph, Leiden
may split locally coherent neigh\-borhoods across multiple communities because
the global modularity gain favors a different partition. DSCF resists this —
the LPA component holds locally coherent groups together even when modularity
would split them.

\textbf{Comparison to LPA-only communities:}

LPA can merge struc\-turally distinct regions if they happen to be locally
connected (the "resolution limit" problem). DSCF resists this — the
modularity signal penalizes over-merging.

\subsubsection{4.3 Probab\-ilistic Traversal: Bayesian Beam Search (v1.1.0)}

To handle topological noise and cold-start uncertainty, CEREBRUM v1.1.0 introduces a proba\-bilistic mode where edge weights are modeled as Beta distr\-ibutions:

\begin{equation}
P(a | \alpha, \beta) = \frac{a^{\alpha-1} (1-a)^{\beta-1}}{B(\alpha, \beta)}
\end{equation}

During traversal, Thompson Sampling is employed to select neighbors. For each candidate neighbor $v$:
\begin{enumerate}
\item A sample $x$ is drawn from $\text{Beta}(\alpha_v, \beta_v)$.
\item Neighbors are ranked by $x$, and the top-$B$ candidates are retained.

\end{enumerate}
\textbf{Warm-Starting}: To reduce initial variance, the distr\-ibution is seeded using the deter\-ministic CSA score $s$:
\begin{equation}
\alpha_{init} = s \cdot \omega, \quad \beta_{init} = (1-s) \cdot \omega
\end{equation}
where $\omega$ is the \texttt{warm\_start\_strength} (default 10.0).

\subsubsection{4.4 Cross-Modal Alignment: Procrustes SVD (v1.1.0)}

\texttt{SignalEncoder} aligns non-textual signals (e.g., FFT spectra of sensor data) into the canonical entity embedding space $\mathcal{E}$ using Orthogonal Procrustes Analysis. Given a set of $n$ anchor points $X$ in the signal space and their corre\-sponding entity embeddings $Y$ in $\mathcal{E}$, we find the optimal rotation matrix $R$:

\begin{equation}
R = \arg\min_{\Omega^T\Omega=I} \| \Omega X - Y \|_F
\end{equation}

The solution is obtained via SVD of the covariance matrix $M = Y X^T$:
\begin{equation}
M = U \Sigma V^T \implies R = U V^T
\end{equation}

\textbf{Canonical Basis Anchor}: To prevent geometric drift in federated hops, all alignments are anchored to a fixed root embedding space $\mathcal{E}_{root}$ rather than chaining through inter\-mediate adapters.

\subsubsection{4.5 Federated Lease \& Pinning (v1.1.0)}

In high-velocity streaming envir\-onments, remote nodes may be evicted before a multi-hop federated query completes. \texttt{FederatedAdapter} implements a lease protocol:
\begin{enumerate}
\item During the discovery phase, the local node sends a \texttt{PIN(entity\_id, ttl)} request to the remote peer.
\item The remote peer "pins" the entity, exempting it from sliding-window eviction for the duration of the TTL.
\item The lease is autom\-atically released upon TTL expiry or explicit \texttt{UNPIN} from the local node.

\end{enumerate}
\subsubsection{4.6 Crypto\-graphic Path Provenance: HMAC-SHA256}

To prevent adversarial path injection in federated reasoning, CEREBRUM implements HMAC-based verif\-ication for remote adapter responses. A shared secret $\mathcal{K}$ is configured between peers. 

For a response body $\mathcal{M}$ (the list of candidate edges), the remote peer computes the signature $\mathcal{S}$:
\begin{equation}
\mathcal{S} = \text{HMAC-SHA256}(\mathcal{K}, \mathcal{M})
\end{equation}

The local adapter verifies the signature by computing $\mathcal{S}' = \text{HMAC-SHA256}(\mathcal{K}, \mathcal{M})$ and checking $\mathcal{S}' \equiv \mathcal{S}$ using a constant-time comparison to prevent timing attacks. Responses with missing or invalid signatures are discarded.

\subsubsection{4.7 Lazy STDP Weight Decay}

The $O(N)$ complexity of global weight decay in the STDP engine is addressed via a lazy update mechanism. For a causal pair $(u, v)$, the system stores the causal weight $w_{uv}$ and the global step counter $t_{last}$ at which the pair was last updated.

Given a current global step $T$ and decay rate $\lambda \in [0, 1]$, the current weight $w'_{uv}$ is computed only upon access:
\begin{equation}
w'_{uv} = w_{uv} \cdot \lambda^{(T - t_{last})}
\end{equation}

This transforms the global $O(N)$ maintenance task into a local $O(1)$ operation, enabling sub-millisecond event processing even on massive, highly-active graphs.

---

\subsection{5. The Forward Pass: Graph Reasoning}

\subsubsection{5.1 Algorithm}

\begin{lstlisting}
INPUT:  Query Q (text string or entity list)
        Graph G (any backend)
        Community assignments C (from DSCF)
        Entity embeddings E (from any KGE method)
        Hop depth L, beam width B, top-K

OUTPUT: Ranked list of reasoning paths P = [(path, score, explanation)]

STEP 1 — Entity Grounding
  If Q is text: extract entities via NER or fuzzy match to graph vocabulary
  S = {e₁, e₂, ..., eₙ}  // seed entities
  h⁰ᵢ = E[eᵢ]             // initial embeddings

STEP 2 — Structural Encoding
  For each seed entity eᵢ:
    pos_i = [pagerank(eᵢ), betweenness(eᵢ), degree(eᵢ), community_id(eᵢ)]
    h⁰ᵢ = LayerNorm(h⁰ᵢ + W_pos · pos_i)

STEP 3 — Attention Traversal (L layers)
  beam = [(path=[eᵢ], embedding=h⁰ᵢ, score=1.0) for each eᵢ in S]

  For k = 1 to L:
    candidates = []
    For each (path, h, score) in beam:
      current = path[-1]
      neighbors = G.neighbors(current)

      For each neighbor v in neighbors:
        w = CSA(current, v, k)                    // attention weight
        h_new = ReLU(W_k · (w · E[v] + h))       // aggregation
        // W_k: hop-depth projection matrix (dim → dim)
        // Zero-shot default: W_k = I (identity) — no learned projection
        // Supervised setting: W_k learned per-hop via path-ranking loss
        h_new += h                                 // residual
        h_new = LayerNorm(h_new)
        path_score = score × w × community_coherence(path + [v])
        candidates.append((path + [v], h_new, path_score))

    beam = top_B(candidates, key=path_score)      // beam pruning

STEP 4 — Path Scoring
  Final score for path P = (e₁ → r₁ → e₂ → r₂ → ... → eₗ):
    score(P) = Π attention_weights
             × community_coherence(P)
             × semantic_alignment(h_final, query_embedding)

STEP 5 — Output
  Return top-K paths ranked by score
  Each path includes:
    - Ordered entity/relation sequence
    - Score breakdown (attention, community, semantic)
    - Community sequence (which "heads" were traversed)
    - Natural language explanation template

**Example Reasoning Trace:**

Query: *"Who discovered a radioactive element?"*

1.  **Entity Grounding**: `Marie Curie` is identified as the seed node.
2.  **Hop 1 (Attention)**: `Marie Curie` **—[discovered]→** `Polonium`.
    - *CSA Weight*: High due to same community membership (Scientific Discoveries) and semantic alignment.
3.  **Hop 2 (Traversal)**: `Polonium` **—[is_a]→** `Radioactive Element`.
    - *CSA Weight*: High due to relation type weights and embedding proximity.
4.  **Output**: The path `Marie Curie → discovered → Polonium → is_a → Radioactive Element` is returned as a grounded proof.
\end{lstlisting}

\subsubsection{5.2 Community Coherence}

The \texttt{community\_coherence} term rewards paths that traverse communities in a
principled way:

\begin{lstlisting}
community_coherence(path) =
    (intra-community steps × 1.0 + cross-community steps × 0.5)
    / total steps
\end{lstlisting}

A path that stays within one community scores 1.0 (tight local reasoning).
A path that makes one community transition scores ~0.75 (one conceptual leap).
This prevents paths that jump incoh\-erently across unrelated domains.

\subsubsection{5.3 Interp\-retability}

The output of CEREBRUM is always a path through the KG. This is funda\-mentally
different from LLM reasoning in two ways:

\begin{enumerate}
\item \textbf{Verifiable}: every step is an explicit (entity, relation, entity) triple
\end{enumerate}
   that exists in the graph. The system cannot hallucinate a connection that
   isn't there.

\begin{enumerate}
\item \textbf{Auditable}: the community sequence tells you \textit{which conceptual domains}
\end{enumerate}
   were traversed. A path that crosses from community "Clinical Trials" to
   community "Drug Mechanisms" to community "Side Effects" is immediately
   under\-standable.

This property makes CEREBRUM speci\-fically valuable in high-stakes domains
(medical, legal, financial) where hallu\-cination is unacc\-eptable.

---

\subsection{6. Implem\-entation Archit\-ecture}

\subsubsection{6.1 Design Principles}

\textbf{Framework agnostic}: CEREBRUM must work with any graph database, any
embedding method, and any LLM (or no LLM). No vendor lock-in.

\textbf{No training required by default}: The zero-shot confi\-guration uses fixed
parameters ($\alpha$, $\beta$, $\gamma$, $\delta$, $\epsilon$) and any entity embedding. Communities are
computed unsup\-ervised via DSCF.

\textbf{Progressive enhancement}: Users can improve performance by providing
training pairs (for learning parameters) or domain-specific edge type weights
without changing the core archi\-tecture.

\textbf{Minimal depen\-dencies}:
\begin{itemize}
\item Core: \texttt{networkx}, \texttt{numpy}, \texttt{leidenalg} (for DSCF), \texttt{scipy}
\item Adapters: optional graph DB drivers
\item Embeddings: optional sentence-trans\-formers or pykeen (for TransE/RotatE)
\item API: optional FastAPI

\end{itemize}
\subsubsection{6.2 Repository Structure}

\begin{lstlisting}
parallax/
│
├── core/
│   ├── __init__.py
│   ├── graph_adapter.py        # Abstract base class for graph backends
│   ├── embedding_engine.py     # Entity embedding interface
│   ├── community_engine.py     # Community detection (DSCF + others)
│   ├── attention_engine.py     # Community-Structured Attention
│   └── structural_encoder.py   # Graph positional encoding
│
├── reasoning/
│   ├── traversal.py            # Beam-search attention traversal
│   ├── path_scorer.py          # Multi-signal path ranking
│   └── answer_extractor.py     # Extract top-K answers
│
├── adapters/
│   ├── networkx_adapter.py     # In-memory (default, no external deps)
│   ├── neo4j_adapter.py        # Neo4j via bolt
│   ├── rdf_adapter.py          # SPARQL endpoint (Wikidata, DBpedia)
│   └── csv_adapter.py          # Bootstrap from edge-list CSV
│
├── llm_bridge/
│   └── context_formatter.py    # Optional: format paths as LLM context
│
├── api/
│   ├── server.py               # FastAPI REST server
│   └── schemas.py              # Pydantic models
│
├── cli/
│   └── parallax.py             # Command-line interface
│
├── tests/
│   ├── test_dscf.py
│   ├── test_csa.py
│   ├── test_traversal.py
│   └── fixtures/toy_graph.csv
│
├── benchmarks/
│   ├── webqsp_eval.py
│   ├── metaqa_eval.py
│   └── baseline_comparison.py
│
├── examples/
│   ├── wikidata_quickstart.py
│   ├── neo4j_quickstart.py
│   └── csv_quickstart.py
│
├── pyproject.toml
├── README.md
└── PAPER.md
\end{lstlisting}

\subsubsection{6.3 The Abstract Graph Adapter}

The adapter pattern is what makes CEREBRUM truly agnostic:

\begin{lstlisting}
from abc import ABC, abstractmethod
from dataclasses import dataclass
from typing import List, Optional

@dataclass
class Entity:
    id: str
    label: str
    type: str
    properties: dict

@dataclass
class Edge:
    source_id: str
    target_id: str
    relation_type: str
    weight: float = 1.0

class GraphAdapter(ABC):
    @abstractmethod
    def get_entity(self, entity_id: str) -> Optional[Entity]: ...

    @abstractmethod
    def get_neighbors(self, entity_id: str,
                      edge_types: List[str] = None,
                      max_neighbors: int = 50) -> List[Edge]: ...

    @abstractmethod
    def find_entities(self, query: str, top_k: int = 10) -> List[Entity]: ...

    @abstractmethod
    def to_networkx(self) -> "nx.Graph": ...  # for community detection
\end{lstlisting}

\subsubsection{6.4 What Comes From Home Assistant}

The following components were originally ported from Home Assistant and generalized:

\begin{itemize}
\item \texttt{tsc\_communities()} — now the Triple-Signal Consensus engine
\item \texttt{\_run\_leiden()} / \texttt{\_run\_lpa()} / \texttt{\_build\_igraph()}
\item Neo4j connection patterns and Cypher templates
\item The community broadcast WebSocket pattern (for live appli\-cations)

\end{itemize}
No other Home Assistant-specific depen\-dencies are carried over. Refer to the \texttt{core/} directory for the production imple\-mentation.

\subsubsection{6.5 Phased Build Plan}

\textbf{Phase 0 — Theory (COMPLETE)}
\begin{itemize}
\item Formalize CSA and TSC; write white paper
\item Prototype consensus logic in Python
\item Validate TSC produces stable communities on Home Assistant's Neo4j graph

\end{itemize}
\textbf{Phase 1 — Core Engine (COMPLETE)}
\begin{itemize}
\item \texttt{core/graph\_adapter.py} — abstract base + NetworkX adapter
\item \texttt{core/community\_engine.py} — TSC, Leiden, LPA, hybrid
\item \texttt{core/embedding\_engine.py} — SentenceEngine (zero-training default)
\item \texttt{core/attention\_engine.py} — CSA weight formula
\item \texttt{core/structural\_encoder.py} — PageRank, betweenness, degree encoding
\item Unit tests on \texttt{fixtures/toy\_graph.csv} for all of the above

\end{itemize}
\textbf{Phase 2 — Reasoning Engine (COMPLETE)}
\begin{itemize}
\item \texttt{reasoning/traversal.py} — beam-search traversal with CSA weights
\item \texttt{reasoning/path\_scorer.py} — multi-signal scoring + community coherence
\item \texttt{reasoning/answer\_extractor.py} — top-K ranked answers
\item Integration test: end-to-end query on toy graph produces grounded paths

\end{itemize}
\textbf{Phase 3 — Adapters + API (COMPLETE)}
\begin{itemize}
\item \texttt{adapters/neo4j\_adapter.py} (port Home Assistant patterns)
\item \texttt{adapters/rdf\_adapter.py} (SPARQL, for Wikidata/DBpedia)
\item \texttt{adapters/csv\_adapter.py} (bootstrap from edge-list)
\item \texttt{api/server.py} — FastAPI REST: \texttt{/query}, \texttt{/communities}, \texttt{/health}
\item \texttt{cli/parallax.py} — command-line interface

\end{itemize}
\textbf{Phase 4 — LLM Bridge + Benchmarks (COMPLETE)}
\begin{itemize}
\item \texttt{llm\_bridge/context\_formatter.py} — format paths as LLM prompts
\item \texttt{benchmarks/webqsp\_eval.py} and \texttt{metaqa\_eval.py}
\item Ablation study: DSCF vs Leiden vs LPA as attention heads
\item Baseline comparisons: BFS, GAT, GraphRAG, vanilla RAG
\item Innovation: Metaedge Bridge Bonus (EF-005) to solve "Type Alignment Trap"

\end{itemize}
\textbf{Phase 5 — Release (COMPLETE)}
\begin{itemize}
\item Final docum\-entation and code cleanup
\item Project-wide validation and stable tag (v0.1.0)
\item Triple-Signal Consensus (TSC) engine rollout

\end{itemize}
\textbf{Phase 6 — Federated Graph Attention (COMPLETE)}
\begin{itemize}
\item \texttt{adapters/federated\_adapter.py} — multi-source aggregation
\item \texttt{adapters/remote\_adapter.py} — HTTP proxy for remote graphs
\item Entity Alignment Index for cross-graph resolution

\end{itemize}
\textbf{Phase 7 — Dynamic Updates (COMPLETE)}
\begin{itemize}
\item Cross-graph "wormhole" attention weights
\item Wildcard community scoring in CSA

\end{itemize}
\textbf{Phase 8 — Holographic Index (COMPLETE)}
\begin{itemize}
\item Bloom Filter based entity probing
\item Community Centroid based semantic discovery
\item Compressed graph signatures for "blind" discovery

\end{itemize}
\textbf{Phase 9 — Stable Release v0.2.0 (COMPLETE)}
\begin{itemize}
\item Handshake protocol for capability negotiation
\item Reasoning Callbacks for cross-graph path verif\-ication
\item Security hardening and full QA audit

\end{itemize}
\subsubsection{6.6 Comput\-ational Complexity}

\textbf{DSCF community detection}: O(E $\cdot$ I) where E = edges, I = iterations
(typically I \textless  50 before convergence). Comparable to Louvain.

\textbf{Community graph const\-ruction}: O(E) post-DSCF; precomputed once.

\textbf{CSA weight per edge}: O(d) where d = embedding dimension. Precompute
structural encodings once; community lookups are O(1) with a hash table.

\textbf{Beam traversal (one query)}: O(B $\cdot$ L $\cdot$ k̄ $\cdot$ d) where:
\begin{itemize}
\item B = beam width (default 10)
\item L = hop depth (default 3)
\item k̄ = average degree
\item d = embedding dimension

\end{itemize}
For a graph with k̄=20, L=3, B=10, d=384: ~230,000 floating-point operations
per query — milli\-seconds on CPU, no GPU required.

\textbf{Comparison to Transformer}: Full attention is O(n² $\cdot$ d) per layer.
CEREBRUM traversal is O(B $\cdot$ L $\cdot$ k̄ $\cdot$ d) — independent of graph size n.
This makes CEREBRUM sublinear in graph size for fixed-width beam search.

\subsubsection{6.7 Known Failure Modes}

\textbf{Dense hub nodes}: Nodes with very high degree (k \textgreater \textgreater  k̄) cause beam
explosion at that hop. Mitigation: cap \texttt{max\_neighbors} per node (default 50).

\textbf{Homogeneous graphs}: If all nodes share the same community, CSA degenerates
to pure embedding similarity (community\_score terms cancel). This occurs in
highly regular graphs (grids, complete bipartite). Mitigation: check community
count post-DSCF; if count \textless  3, fall back to BFS with embedding similarity only.

\textbf{Discon\-nected graphs}: Multiple connected components each get their own
communities; cross-component traversal is impossible by definition. Queries
spanning components will return no paths. Mitigation: surface component
boundaries to callers; allow separate per-component queries.

\textbf{Sparse embedding coverage}: Entities with generic or missing labels get
near-random sentence embeddings. The $\alpha$ term (embedding similarity) becomes
noise; DSCF community structure ($\beta$ term) carries the full weight. Still
functions; inter\-pretability of similarity scores is reduced.

\textbf{Adversarial community injection}: A malicious actor who can insert edges
into the graph can manipulate community structure and thus attention weights.
Relevant for appli\-cations where the KG is user-writable. Mitigation: sign
trusted edges; treat unsigned-edge-induced communities with lower $\beta$ weight.

\subsubsection{6.8 Horizontal Scalability}

\textbf{BLUF: CEREBRUM query serving scales horiz\-ontally without archi\-tectural limit. Queries are stateless reads from a sharded graph database; throughput grows linearly with workers. LLM inference cannot do this.}

CEREBRUM's scalability is bounded by the scalability of graph databases — a solved problem at web scale — not by the constraints of neural network inference.

\paragraph{Three independent parallelism dimensions}

\textbf{Within a single query.} At each beam hop, all B$\times$k̄ candidate edge scores are computed indep\-endently. The only synch\-ronization point is the top-B selection (O(B log B) on ~200 items). Every other operation — embedding lookup, cosine similarity, CSA weight computation — is embar\-rassingly parallel across GPU threads.

\begin{lstlisting}
Hop k:
  Score(candidate_1 × neighbors)  ──┐
  Score(candidate_2 × neighbors)  ──┤→ top-B select → next hop
  ...                               │
  Score(candidate_B × neighbors)  ──┘
  (only sync point: O(B log B) reduce)
\end{lstlisting}

\textbf{Across concurrent queries.} Traversal workers are stateless during inference — the graph and embedding matrix are read-only. Multiple concurrent queries share no mutable state. Query throughput scales linearly with worker count up to I/O bandwidth limits, with zero inter-query synch\-ronization overhead.

\textbf{Across machines.} The CSA formula strongly prefers staying within a community (community\_score = 1.0 intra vs. exponential decay inter). This makes communities the natural shard key: most queries never cross a shard boundary and therefore never leave their home machine. Cross-shard traversal is a sparse, predictable network call — already implemented in the FederatedAdapter. The holographic index (Bloom filter + community centroid) lets any machine identify which remote shard owns an entity without loading remote data.

\paragraph{Amdahl's Law comparison}

\begin{center}
{\footnotesize
\begin{tabularx}{\columnwidth}{|>{\raggedright\arraybackslash}X|>{\raggedright\arraybackslash}X|>{\raggedright\arraybackslash}X|}
\hline
 & LLM inference & CEREBRUM query \\ \hline
Sequential fraction & ~100\% of forward pass (layer N waits for layer N-1) & Top-B select at each of L hops (~3$\mu$s total) \\ \hline
Parallel fraction & ~0\% within a single inference call & ~99.9\% of all scoring work \\ \hline
Cross-machine sync & All-reduce at every layer boundary & Sparse shard crossings (most queries: zero) \\ \hline
Practical speedup ceiling & ~8–16 GPUs per inference call (commu\-nication dominates) & Linear with worker count up to DB I/O limits \\ \hline
\end{tabularx}
}
\end{center}

\paragraph{Scaling tiers}

\begin{center}
{\footnotesize
\begin{tabularx}{\columnwidth}{|>{\raggedright\arraybackslash}X|>{\raggedright\arraybackslash}X|>{\raggedright\arraybackslash}X|>{\raggedright\arraybackslash}X|}
\hline
Scale & Graph storage & Compute & Status \\ \hline
Millions of nodes & NetworkXAdapter (RAM) & Single machine, GPU embeddings & Implemented \\ \hline
Billions of edges & Neo4jAdapter (disk) & Single machine & Adapter exists \\ \hline
Multi-machine partitioned & FederatedAdapter (community-sharded) & Cluster of API instances & Phase 6–9 \\ \hline
100B+ edges & Distributed graph DB (Neptune, TigerGraph) + stateless Kubernetes workers & Cluster & Archit\-ecture compatible; adapter needed \\ \hline
Web scale (1T+ edges) & Distributed graph processing (Spark GraphX) for offline DSCF; standard graph DB for queries & Datacenter & Theoretical; DSCF distr\-ibution is the hard problem \\ \hline
\end{tabularx}
}
\end{center}

The bottleneck at extreme scale shifts entirely to community detection — DSCF requires a global view of the graph to compute modularity gain. Distributed community detection (distributed Louvain, FENNEL-style parti\-tioning) is an active research area. Critically, community detection runs \textbf{offline and perio\-dically} — queries continue running against the last known partition while the next detection pass runs in the background. It is never on the query critical path.

\textbf{The one-line summary: a KG query touches ~0.001\% of the graph per call using embar\-rassingly parallel operations. An LLM touches 100\% of its parameters every time, seque\-ntially. That gap defines the scalability story.}

---

\subsection{7. Experi\-mental Results}

\subsubsection{7.0 Experi\-mental Environment}

All benchmarks were executed on the following hardware and software confi\-guration to support repro\-ducibility:

\begin{itemize}
\item \textbf{CPU}: AMD Ryzen 9 9950X3D 16-Core Processor (32 Logical Processors)
\item \textbf{RAM}: 64 GB DDR5
\item \textbf{OS}: Windows 11 Pro (Build 10.0.26220)
\item \textbf{Python}: 3.14.0
\item \textbf{Graph Backends}: NetworkX 3.4.2, igraph 0.11.6
\item \textbf{Embeddings}: RandomEngine (64-dim) for structural validation; SentenceEngine (384-dim) for semantic tasks.

\end{itemize}
\subsubsection{7.1 MetaQA: The Baseline Lower Bound}

MetaQA evaluation revealed a \textbf{Structural Mismatch (EF-004)}. Because MetaQA answer paths always cross entity-type boundaries (Movie $\to$ Actor), and community detection naturally separates these types, the default CSA formula (favoring intra-community edges) penalized the correct paths.

\begin{itemize}
\item \textbf{Outcome}: BFS outpe\-rformed CSA variants on Hits@1.
\item \textbf{Signif\-icance}: Established the "lower bound" of performance on topologies where community signal is anti-informative.

\end{itemize}
\subsubsection{7.2 The Bridge Bonus Innovation (EF-005)}

To solve the "Type Alignment Trap" identified in MetaQA and Hetionet, we introduced the \textbf{Metaedge Bridge Bonus} ($w_{rel}$ in the CSA formula). By assigning a positive bonus (e.g., 0.4) to inter-type metaedges like \texttt{treats} or \texttt{associates}, we offset the cross-community penalty while retaining structural guidance.

\subsubsection{7.3 Hetionet: Biomedical Reasoning at Scale}

On a 500,000-edge subset of Hetionet, CEREBRUM with LPA attention heads and the Bridge Bonus signi\-ficantly outpe\-rformed the BFS baseline.

\begin{center}
{\footnotesize
\begin{tabularx}{\columnwidth}{|>{\raggedright\arraybackslash}X|l|>{\raggedright\arraybackslash}X|>{\raggedright\arraybackslash}X|>{\raggedright\arraybackslash}X|}
\hline
Template & Hop & LPA+CSA H@1 & BFS H@1 & Improvement \\ \hline
disease\_associates\_gene & 1 & \textbf{0.6560} & 0.4320 & \textbf{+51.8\%} \\ \hline
gene\_participates\_pathway & 1 & \textbf{0.2600} & 0.0950 & \textbf{+173.6\%} \\ \hline
\end{tabularx}
}
\end{center}

\subsubsection{7.4 WebQSP: Real-world Entity Lookup}

On the WebQSP benchmark (FB15k-237), CEREBRUM demon\-strated superior recall and ranking quality.

\begin{center}
{\footnotesize
\begin{tabularx}{\columnwidth}{|>{\raggedright\arraybackslash}X|>{\raggedright\arraybackslash}X|>{\raggedright\arraybackslash}X|>{\raggedright\arraybackslash}X|}
\hline
Variant & Hits@1 & Hits@10 & MRR \\ \hline
CEREBRUM (LPA+CSA) & 0.0400 & \textbf{0.3360} & \textbf{0.1203} \\ \hline
BFS Baseline & 0.0540 & 0.3000 & 0.1081 \\ \hline
\end{tabularx}
}
\end{center}

\subsubsection{7.5 Key Findings}

\begin{enumerate}
\item \textbf{Recall Advantage}: CSA variants consi\-stently achieve higher recall (Hits@10) than BFS, validating the system's ability to steer the beam toward correct graph regions.
\item \textbf{Signal Duality}: DSCF provides finer-grained precision, while LPA provides coarser, more robust recall.
\item \textbf{Zero-Shot Viability}: All results were achieved using random embeddings and manual weights, proving CEREBRUM works without any training data.

\end{enumerate}
---

\subsection{8. The DSCF-as-Attention-Head Hypothesis}

The central theoretical claim of CEREBRUM is that DSCF communities are better
attention heads than Leiden-only or LPA-only communities. We state this as a
falsifiable hypothesis:

\textbf{H1 (DSCF Attention Hypothesis)}: For multi-hop reasoning tasks on KGs,
CEREBRUM with DSCF attention heads achieves higher answer accuracy than
CEREBRUM with Leiden-only or LPA-only attention heads.

\textbf{H2 (CSA vs GAT Hypothesis)}: CSA-guided traversal achieves higher accuracy
on multi-hop questions than GAT-based traversal on the same graph and same
entity embeddings.

\textbf{H3 (Interp\-retability Hypothesis)}: CEREBRUM paths receive higher human
coherence ratings than equivalent LLM-generated reasoning chains on the same
questions, because every step is a grounded graph edge.

\textbf{H3 Evaluation Protocol}: Present matched pairs — one CEREBRUM path and
one LLM reasoning chain — to N$\ge$30 annotators blind to their source. Ask:
"Which reasoning chain is more coherent and trustworthy? (A / B / Equal)".
Primary metric: proportion preferring CEREBRUM path. Secondary: Cohen's kappa
for inter-annotator agreement (target $\kappa$ \textgreater  0.6).

These hypotheses are testable on standard benchmarks (WebQSP, MetaQA-3hop)
and define the empirical work for Phase 2.

---

\subsection{9. Benchmarks and Results}

\subsubsection{9.1 Datasets}

\begin{center}
{\footnotesize
\begin{tabularx}{\columnwidth}{|>{\raggedright\arraybackslash}X|>{\raggedright\arraybackslash}X|l|>{\raggedright\arraybackslash}X|}
\hline
Dataset & Task & Hops & Size \\ \hline
WebQSP & Single + multi-hop QA & 1-2 & 4,737 questions \\ \hline
MetaQA-2hop & Multi-hop QA & 2 & 118,980 questions \\ \hline
MetaQA-3hop & Multi-hop QA & 3 & 114,196 questions \\ \hline
FB15k-237 & Link prediction & - & 310,116 triples \\ \hline
Toy graph (internal) & Unit testing & 1-4 & ~200 nodes \\ \hline
\end{tabularx}
}
\end{center}

\subsubsection{9.2 Baselines}

\begin{center}
{\footnotesize
\begin{tabularx}{\columnwidth}{|>{\raggedright\arraybackslash}X|>{\raggedright\arraybackslash}X|>{\raggedright\arraybackslash}X|}
\hline
System & Type & Notes \\ \hline
BFS (no attention) & Graph traversal & Traversal without CSA weighting \\ \hline
GAT & Graph neural network & 2-layer, trained \\ \hline
GraphRAG & LLM-based & Community summaries $\to$ GPT-4 \\ \hline
RAG (vanilla) & LLM-based & FAISS retrieval $\to$ GPT-4 \\ \hline
\textbf{CEREBRUM (TSC)} & Graph attention & Ours, Triple-Signal Consensus heads \\ \hline
\textbf{CEREBRUM (LPA)} & Graph attention & Ablation: LPA-only heads \\ \hline
\end{tabularx}
}
\end{center}

\subsubsection{9.3 Metrics}

\begin{itemize}
\item \textbf{Hits@1, Hits@3, Hits@10}: answer in top-K paths
\item \textbf{Mean Reciprocal Rank (MRR)}: ranked answer quality
\item \textbf{Path coherence} (human eval): are the reasoning paths under\-standable?
\item \textbf{Grounding rate}: what fraction of returned paths are fully grounded
\end{itemize}
  (all edges verified in the KG)?

\subsubsection{9.4 Phase 4 Results (Ablation Study)}

A rigorous ablation study on MetaQA (Run 019) yielded the following key engineering findings:

\begin{enumerate}
\item \textbf{Structural Mismatch (EF-004)}: Breadth-First Search (BFS) consi\-stently outperforms CSA variants on MetaQA. This confirms that MetaQA's question structure (cross-type entity lookup) is penalized by community-based attention, which favors intra-community coherence. This is a dataset-specific chara\-cteristic, not an algorithmic defect.
\item \textbf{TSC Stability}: The Triple-Signal Consensus (TSC) engine demon\-strated superior stability in community detection compared to earlier DSCF iterations, producing consistent partition counts across runs.
\item \textbf{The Mesoscale Gap}: TSC produces fine-grained communities (~14k on MetaQA) compared to LPA (~1.6k). This granularity provides high precision for local queries but neces\-sitates the "Metaedge Bridge Bonus" or Federated Reasoning strategies to bridge large topological distances in multi-hop tasks.

\end{enumerate}
---

\subsection{10. Open Research Questions}

\subsubsection{9.1 Embedding Strategy}

Two options exist:

\textbf{Option A — Pre-trained structural embeddings (TransE/RotatE)}: trained
on the graph structure itself. More precise but requires a training step.
Suitable for static KGs or when training compute is available.

\textbf{Option B — On-the-fly label embeddings (sentence-trans\-formers)}: encode
entity labels and descr\-iptions using a pre-trained language model. No graph-
specific training needed. More agnostic. Less precise for entities with
ambiguous labels.

\textbf{Recommended default}: Option B for zero-shot deployment; Option A when the
graph has been stable and training is feasible. CEREBRUM should support both
inter\-changeably via the EmbeddingEngine interface.

\subsubsection{8.2 Adaptive Community Granularity}

The DSCF resolution parameter controls how many communities are formed. Too
few communities = coarse attention heads that miss structure. Too many = noisy
heads that don't generalize.

\textbf{Proposed adaptive rule}: target K $\approx$ √N communities, where N = node count.
This is consistent with theoretical results on optimal modularity resolution
and supports attention head count scales sensibly with graph size.

For Home Assistant's KG (N $\approx$ 5,000): target ~70 communities.
For Wikidata subset (N $\approx$ 100,000): target ~316 communities.

\subsubsection{8.3 Soft vs Hard Community Membership}

DSCF produces hard assignments (each node belongs to exactly one community).
Real-world entities often span multiple communities — a person can be both
a scientist and a politician.

\textbf{Extension}: weight-based soft membership, where each node has a probability
distr\-ibution over communities. The community\_score function becomes a
dot product of membership vectors. This would require modifying DSCF to track
confidence scores at convergence.

\subsubsection{8.4 Learnable Parameters}

In zero-shot mode, $\alpha$, $\beta$, $\gamma$, $\delta$, $\epsilon$ are fixed. For supervised settings, they
can be learned from (query, ground-truth-answer) pairs via gradient descent
on a path-ranking loss. This is an optional enhancement that does not affect
the core archi\-tecture.

\subsubsection{8.5 Temporal Knowledge Graphs}

Time-stamped KGs (events, evolving relat\-ionships) introduce a temporal
dimension. The positional encoding would need to incorporate temporal distance
alongside graph-structural distance. Left for future work.

---

\subsection{Acknow\-ledgments: Intell\-ectual Debt and Credits}

CEREBRUM stands on the shoulders of decades of research in graph theory, community detection, and neural networks. We explicitly acknowledge the found\-ational work of the following researchers and the algorithms that form the bedrock of our framework:

\begin{enumerate}
\item \textbf{LPA (Label Propagation Algorithm)}: Usha Nandini Raghavan, Réka Albert, and Shailesh Kumara (2007). Their work on near-linear time community detection via local neighbor voting provided the "Local Signal" for our DSCF engine.
\item \textbf{Louvain Algorithm}: Vincent Blondel, Jean-Loup Guillaume, Renaud Lambiotte, and Etienne Lefebvre (2008). Their greedy modularity optim\-ization method established the global structural baseline for community detection.
\item \textbf{Leiden Algorithm}: Vincent Traag, Ludo Waltman, and Nees Jan van Eck (2019). Their refinement of Louvain, ensuring internal conne\-ctivity, provides the "Global Signal" and conne\-ctivity post-pass for DSCF.
\item \textbf{Graph Attention Networks (GATs)}: Petar Veličković, Guillem Cucurull, Arantxa Casanova, Adriana Romero, Pietro Liò, and Yoshua Bengio (2018). Their intro\-duction of learned attention on graphs served as the primary foil and inspiration for our Community-Structured Attention (CSA).
\item \textbf{KG Embeddings (TransE / RotatE)}: Antoine Bordes et al. (2013) and Zhiqing Sun et al. (2019). Their work on repre\-senting relational knowledge in vector spaces provides the semantic grounding layer for CSA.
\item \textbf{GraphRAG}: Microsoft Research / Edge et al. (2024). Their pioneering work in combining community summaries with LLM retrieval provided the immediate context and competitive baseline for CEREBRUM's grounded reasoning approach.
\item \textbf{Avionics Engineering}: The concept of \textbf{mid-level voting} (or mid-value selection) in triplex-redundant aircraft navigation. This engineering principle of multi-sensor consensus served as the found\-ational inspiration for the multi-signal logic of DSCF and TSC.
\item \textbf{PageRank}: Page, L., Brin, S., Motwani, R., \& Winograd, T. (1999). The PageRank Citation Ranking: Bringing Order to the Web. Stanford InfoLab. Used as the global authority prior in the CSA formula and as a THALAMUS structural encoding feature.
\item \textbf{Betweenness Centrality}: Freeman, L.C. (1977). A set of measures of centrality based on betweenness. \textit{Sociometry}, 40(1), 35–41. Used as a THALAMUS positional encoding feature alongside PageRank.
\item \textbf{Simulated Annealing}: Kirkpatrick, S., Gelatt, C.D., \& Vecchi, M.P. (1983). Optimi\-zation by simulated annealing. \textit{Science}, 220(4598), 671–680. The DSCF temperature decay schedule ($\tau$ $\times$ 0.92 per iteration) is a direct application.
\item \textbf{Bloom Filters}: Bloom, B.H. (1970). Space/time trade-offs in hash coding with allowable errors. \textit{Commun\-ications of the ACM}, 13(7), 422–426. Used in Hologr\-aphicIndex for federated blind graph discovery.
\item \textbf{Spike-Timing Dependent Plasticity (STDP)}: Bi, G.Q. \& Poo, M.M. (1998). Synaptic modif\-ications in cultured hippocampal neurons. \textit{Journal of Neuros\-cience}, 18(24), 10464–10472. Markram, H. et al. (1997). Regulation of synaptic efficacy by coincidence of posts\-ynaptic APs and EPSPs. \textit{Science}, 275(5297), 213–215. The STDPDiscre\-tizer and Bridge Twin formation are direct compu\-tational analogs.
\item \textbf{Hebbian Learning}: Hebb, D.O. (1949). \textit{The Organi\-zation of Behavior}. Wiley. The Bridge Twin LTP/LTD analog and InsightEngine Hebbian reward propagation are grounded in this principle.
\item \textbf{Beam Search}: Lowerre, B.T. (1976). The HARPY Speech Recognition System. PhD thesis, Carnegie Mellon University. BeamTraversal is a direct application of this classical algorithm to graph path expansion.
\item \textbf{Sentence-BERT}: Reimers, N. \& Gurevych, I. (2019). Sentence-BERT: Sentence Embeddings using Siamese BERT-Networks. \textit{EMNLP 2019}. Used as the default embedding backend when the \texttt{[embeddings]} extra is installed.


\end{enumerate}
---

\subsection{8.6 Triple-Signal Consensus (TSC): The Next Frontier}

A significant archi\-tectural expansion for CEREBRUM is the transition from the dual-signal DSCF to a \textbf{Triple-Signal Consensus (TSC)} framework. This evolution is designed to close the \textbf{"Mesoscale Gap"}—the structural region between immediate local topology (LPA) and global modularity (Leiden).

\subsubsection{The Motivation for a Third Signal}

While modularity (Global) captures static edge density and LPA (Local) captures immediate neigh\-borhood cohesion, they both miss the \textbf{dynamic flow of information} through a network. In reasoning tasks, the most relevant path is often the one that information naturally "flows" along.

\subsubsection{The TSC Components}

\begin{enumerate}
\item \textbf{LPA (Local)}: Neighbor recognition (Cohesion).
\item \textbf{Modularity (Global)}: Archit\-ecture optim\-ization (Signif\-icance).
\item \textbf{Infomap / Map Equation (Mid-Level)}: Flow-based clustering (Connec\-tivity). Originally proposed by Martin Rosvall and Carl Bergstrom (2008), Infomap uses random walks to identify sub-clusters based on information flow, acting as the "mesoscale" judge between local and global signals.

\end{enumerate}
\subsubsection{Consensus Decision Logic}

In the TSC framework, a node move or a traversal edge must pass a \textbf{Consensus Filter}. This reduces "structural hallu\-cinations"—paths that exist topol\-ogically but lack conceptual or infor\-mational flow coherence.

The fused probability for a node move or attention weight calculation becomes:
\begin{equation}
P(\text{move}) = f(\text{LPA} \cdot \tau_{local}, \text{Mod} \cdot \tau_{global}, \text{Infomap} \cdot \tau_{mid})
\end{equation}

This "mid-level voting" helps confirm that only the most struc\-turally and dynamically robust reasoning chains survive the beam-search pruning process. TSC will be implemented as an optional, high-precision mode within the CEREBRUM core, allowing for direct comparison with DSCF.

\subsubsection{TSC Development Roadmap (Phases 6–10)}

\begin{center}
{\footnotesize
\begin{tabularx}{\columnwidth}{|>{\raggedright\arraybackslash}X|>{\raggedright\arraybackslash}X|>{\raggedright\arraybackslash}X|}
\hline
Phase & Title & Primary Research Objective \\ \hline
\textbf{Phase 6} & \textbf{Flow Integration} & Formalize $P(\text{move}) = f(\text{LPA, Mod, Infomap})$ decision logic. \\ \hline
\textbf{Phase 7} & \textbf{TSC Engine} & Implement Infomap signal processing and Consensus Confidence scoring. \\ \hline
\textbf{Phase 8} & \textbf{Dynamic Traversal} & Implement consensus-gated Beam Search and \texttt{flow\_coherence} ranking. \\ \hline
\textbf{Phase 9} & \textbf{Ablation H4} & Prove precision superiority of TSC over DSCF on complex topologies. \\ \hline
\textbf{Phase 10} & \textbf{v0.2.0 Release} & Produc\-tionize high-precision reasoning for high-stakes deployment. \\ \hline
\end{tabularx}
}
\end{center}

---

---

\subsection{10. Broader Impact and Applic\-ations}

\subsubsection{10.1 Domain Applic\-ations}

\textbf{Biomedical}: Drug-gene-disease-pathway graphs. Multi-hop reasoning for drug
repurposing ("Drug X inhibits enzyme Y which is overe\-xpressed in disease Z").
Grounded inference is critical — no LLM should hallucinate drug inter\-actions.

\textbf{Legal}: Case law citation and statutory reference networks. Multi-hop
precedent tracing. Every step in a legal argument must be citable; CEREBRUM's
grounded paths match this requirement exactly.

\textbf{Cybers\-ecurity}: Attack graphs, CVE dependency networks. "What path leads
from this exposed service to root access?" — a life-safety question that
benefits from verified, traceable reasoning chains.

\textbf{Software engineering}: Code dependency and call graphs. Impact analysis:
"What does changing function X affect?" traversed as a multi-hop attention
path with community context (same module = high attention).

\textbf{Finance}: Entity relat\-ionship graphs for regulatory compliance. Traceable
reasoning chains for auditors: "Why did this transaction trigger a flag?"

\subsubsection{10.2 The LLM Bridge}

CEREBRUM is designed to augment LLMs, not replace them. The \texttt{llm\_bridge}
module formats traversal output as structured context:

\begin{lstlisting}
You are reasoning about: [query]

The knowledge graph traversal found these paths:

Path 1 (score: 0.94):
  Marie Curie [COMMUNITY: Scientific Discoveries]
  → [discovered] →
  Polonium [COMMUNITY: Scientific Discoveries]
  → [exhibits] →
  Radioactivity [COMMUNITY: Physics Phenomena]

Please summarize what this tells us about [query] in natural language.
\end{lstlisting}

This gives any LLM a grounded, structured context that minimizes the risk of
hallu\-cination because the facts are provided explicitly. The LLM's role is
purely natural language generation, not reasoning.

\subsubsection{10.3 The Agnosticism Property}

CEREBRUM is agnostic across five dimensions:

\begin{enumerate}
\item \textbf{Graph database}: implement GraphAdapter for any system
\item \textbf{Embedding method}: implement EmbeddingEngine for any model
\item \textbf{LLM}: any model or none — CEREBRUM works without one
\item \textbf{Domain}: the algorithm is domain-blind; community structure emerges
\end{enumerate}
   from the graph's own topology
\begin{enumerate}
\item \textbf{Query language}: entities can be identified from text, IDs, or direct
\end{enumerate}
   lookup — the entry point is flexible

---

\subsection{11. Production Hardening (Phase 10)}

Phase 10 trans\-itioned CEREBRUM from research prototype to deployable service.

\subsubsection{11.1 JWT Authen\-tication}

All REST endpoints require a Bearer JWT. Tokens are validated by a FastAPI
dependency; expired tokens return HTTP 403, missing tokens return HTTP 401.
Token issuance is external, enabling OAuth2/OIDC integration.

\subsubsection{11.2 ResourceGovernor}

Per-request resource ceilings prevent query runaway:

\begin{equation}
\text{ResourceGovernor: } \{N_{\max}, E_{\max}, B_{\max}, H_{\max}, T_{\max}, C_{\max}\}
\end{equation}

where $N_{\max}$ = max nodes, $E_{\max}$ = max edges, $B_{\max}$ = beam width,
$H_{\max}$ = max hops, $T_{\max}$ = timeout (seconds), $C_{\max}$ = concurrent
queries. Violations return HTTP 429 with a structured JSON error body.

\subsubsection{11.3 Asynch\-ronous Beam Traversal}

\texttt{AsyncBeamTraversal} wraps \texttt{BeamTraversal} in \texttt{asyncio} and adds a
\texttt{/query/stream} SSE endpoint. One event is emitted per beam step:

\begin{lstlisting}
event: step
data: {"hop": k, "candidates": [...], "scores": [...]}

event: result
data: {"paths": [...], "top_score": 0.94, "hops": 2}
\end{lstlisting}

The \texttt{ResourceGovernor} timeout applies per stream connection.

---

\subsection{12. Real-Time Streaming (Phase 11)}

Phase 11 extends CEREBRUM to continuous, live data. The streaming subsystem
sits above the core CSA/DSCF algorithms, which are unchanged.

\subsubsection{12.1 Archit\-ecture Layers}

\begin{center}
{\footnotesize
\begin{tabularx}{\columnwidth}{|>{\raggedright\arraybackslash}X|>{\raggedright\arraybackslash}X|>{\raggedright\arraybackslash}X|}
\hline
Layer & Component & Respon\-sibility \\ \hline
Ingestion & \texttt{StreamSource} subclasses & FileTail, HTTP polling, WebSocket, MQTT, Python callbacks \\ \hline
Discre\-tization & \texttt{*Discretizer} classes & Continuous signal $\to$ discrete graph edges \\ \hline
Buffering & \texttt{SlidingWindowBuffer} & Time-window + count cap + reference counting \\ \hline
Live graph & \texttt{StreamAdapter} & Thread-safe \texttt{NetworkXAdapter} with mutation broadcast \\ \hline
Community & \texttt{IncrementalCommunityUpdater} & Ego-network DSCF on affected nodes only \\ \hline
API & \texttt{/stream/*} endpoints & Ingest, status, SSE event subsc\-ription \\ \hline
\end{tabularx}
}
\end{center}

\subsubsection{12.2 Sliding Window Buffer}

The buffer maintains live edges satisfying two simul\-taneous constraints:

\begin{equation}
\mathcal{W}(t) = \left\{ e \in \mathcal{E} \;\middle|\; t - t_e \leq \Delta t \right\} \cap \left\{ |\mathcal{E}| \leq N_{\max} \right\}
\end{equation}

Reference counting allows the same edge $(s, r, t)$ to be held by multiple
overlapping events; the edge is evicted only when all references expire.

\subsubsection{12.3 Signal Discre\-tization}

\textbf{ThresholdDiscre\-tizer} — categorical states with hysteresis:

\begin{equation}
\text{state}(x) = \begin{cases} \text{SPIKE} & x > \mu + 3\sigma \\ \text{HIGH} & x > \mu + \sigma \\ \text{LOW} & x < \mu - \sigma \\ \text{NORMAL} & \text{otherwise} \end{cases}
\end{equation}

\textbf{BinningDiscre\-tizer} — uniform quant\-ization into $N$ bins:

\begin{equation}
\text{bin}(x) = \left\lfloor \frac{x - x_{\min}}{x_{\max} - x_{\min}} \cdot N \right\rfloor
\end{equation}

\textbf{CoActivationDiscre\-tizer} — emit \texttt{CO\_ACTIVATES} when two sensors fire within
window $\delta_t$ and have co-activated at least $n_{\min}$ times:

\begin{equation}
\text{emit}(A,B) \iff |t_A - t_B| \leq \delta_t \;\wedge\; \text{count}(A,B) \geq n_{\min}
\end{equation}

\textbf{TemporalSequenceDiscre\-tizer} — \texttt{PRECEDES} chains between consecutive events.

\textbf{ObjectDetectionDiscre\-tizer} — \texttt{DETECTS} + \texttt{CO\_OCCURS\_WITH} from frame detections.

\subsubsection{12.4 Incremental Community Updates}

Full DSCF after each ingestion event is prohibitive. The updater restricts
re-computation to the ego-network of radius $r$ around affected nodes:

\begin{equation}
G_{\text{local}} = \{v \mid d_G(v, u) \leq r,\; u \in \mathcal{A}\}
\end{equation}

Updates trigger when $|\mathcal{A}| \geq n_{\min}$ (default 10). All
assignments outside $G_{\text{local}}$ are preserved from the previous full
run, maintaining a globally-consistent, locally-fresh partition.

\subsubsection{12.5 Streaming API}

\begin{center}
{\footnotesize
\begin{tabularx}{\columnwidth}{|>{\raggedright\arraybackslash}X|>{\raggedright\arraybackslash}X|>{\raggedright\arraybackslash}X|}
\hline
Endpoint & Method & Description \\ \hline
\texttt{/stream/ingest} & POST & Batch-ingest \texttt{StreamEvent} objects \\ \hline
\texttt{/stream/status} & GET & Live stats (nodes, edges, events/s, communities) \\ \hline
\texttt{/stream/events} & GET (SSE) & Subscribe to graph mutation events \\ \hline
\end{tabularx}
}
\end{center}

---

\subsection{12.6 Bridge Twin Nodes (Phase 12)}

When a node crosses community boundaries repeatedly — because the graph's true reasoning path requires it — the CSA distance penalty imposes a persistent cost on an edge that is struc\-turally correct. \textbf{Bridge Twins} resolve this by mater\-ialising a copy of the node inside the destination community once the crossing count and semantic fit threshold are both met:

\begin{equation}
\text{fit}(v, \mathcal{C}_d) = \cos\!\left(\mathbf{e}_v,\; \frac{1}{|\mathcal{C}_d|}\sum_{u \in \mathcal{C}_d} \mathbf{e}_u\right) \geq \theta_{bridge}
\end{equation}

The twin $v'$ is assigned $c(v') = \mathcal{C}_d$, carries $\mathbf{e}_{v'} = \mathbf{e}_v$, and is connected via bidir\-ectional \texttt{BRIDGE\_TWIN} edges ($w_{rel} = 1.0$). The CSA short-circuits to $\sigma(0.925 + \varepsilon\phi(k)) \approx 0.716$ for bridge edges. Idle bridges are pruned after $\tau_{prune}$ days (LTD analog). Biological corre\-spondence: thalamic relay nuclei (LGN).

---

\subsection{12.7 Spike-Timing Dependent Plasticity (Phase 13)}

The \texttt{CoActivationDiscretizer} detects symmetric corre\-lations. \texttt{STDPDiscretizer} adds \textbf{directional causal inference} from timing alone, matching the biological STDP rule:

$$\Delta w(A \to B) = \begin{cases}
A\_+ \exp(-\Delta t/\tau\_+) \& \Delta t \textgreater  0 \quad \text{(LTP: A precedes B)} \\
-A\_- \exp(\Delta t/\tau\_-) \& \Delta t \leq 0 \quad \text{(LTD: anti-causal)}
\end{cases}$$

Default: $A_+ = 0.1$, $A_- = 0.105$ (slight LTD dominance, matching Bi \& Poo 1998). A \texttt{CAUSES(A→B)} edge is emitted when $w(A \to B) \geq w_{threshold}$ AND $n(A \to B) \geq n_{min}$. Per-spike multi\-plicative weight decay ($\lambda \in (0,1]$) provides exponential forgetting. The result: CEREBRUM auton\-omously infers directed causal chains from streaming data without domain confi\-guration or labeled training.

---

\subsection{12.8 REM Cycle Memory Consol\-idation (Phase 14)}

The \texttt{REMEngine} (\texttt{core/rem\_engine.py}) implements a biolo\-gically-inspired offline memory conso\-lidation pass modeled on NREM slow-wave sleep. It operates in three sequential phases: \textbf{Prune}, \textbf{Consolidate}, and \textbf{Synthesize}.

\textbf{Prune} removes edges whose confidence has decayed below a floor threshold:

\begin{equation}
\text{prune}(e) \iff \text{confidence}(e) < \theta_{prune} = 0.2 \;\wedge\; \text{type}(e) \neq \text{BRIDGE\_TWIN}
\end{equation}

\texttt{BRIDGE\_TWIN} edges are immune — they represent experience-confirmed structural relays (Phase 12) and must survive the prune pass.

\textbf{Consolidate} re-runs DSCF on the pruned graph so that community boundaries re-align with the surviving, higher-confidence edge set.

\textbf{Synthesize} scans all surviving entity pairs and emits a \texttt{rem\_synthesized} edge when two conditions are jointly met:

\begin{equation}
\cos(\mathbf{e}_u, \mathbf{e}_v) \geq \theta_{sim} = 0.8 \;\wedge\; d_G(u, v) \leq h_{prox} = 4
\end{equation}

Synthesized edges carry \texttt{confidence = 0.3}, placing them just above the prune floor and making them candidates for validation by subsequent traversal.

\textbf{Operational modes:}

\begin{center}
\begin{tabularx}{\columnwidth}{|>{\raggedright\arraybackslash}X|>{\raggedright\arraybackslash}X|}
\hline
Mode & Effect \\ \hline
\texttt{dry\_run=True} & Zero mutations; returns full report of what would be pruned, re-conso\-lidated, and synthesized \\ \hline
\texttt{dry\_run=False} & Mutations applied in-place \\ \hline
\texttt{rollback()} & Restores pruned edges with original attributes; removes synthetic edges. One level of undo. \\ \hline
\end{tabularx}
\end{center}

\textbf{API endpoints:}

\begin{center}
{\footnotesize
\begin{tabularx}{\columnwidth}{|>{\raggedright\arraybackslash}X|>{\raggedright\arraybackslash}X|>{\raggedright\arraybackslash}X|}
\hline
Endpoint & Method & Description \\ \hline
\texttt{/rem/run} & POST & Execute REM cycle (\texttt{dry\_run} param, default \texttt{true}) \\ \hline
\texttt{/rem/rollback} & POST & Undo last non-dry-run cycle \\ \hline
\texttt{/rem/status} & GET & Current REM state, last cycle stats \\ \hline
\end{tabularx}
}
\end{center}

Biological analog: \textbf{NREM slow-wave sleep} — prune weak synapses (synaptic homeostasis), consolidate the community archi\-tecture, and form associative links between seman\-tically proximate but topol\-ogically distant nodes.

---

\subsection{12.9 InsightEngine — Surprise-Driven Discovery (Phase 15)}

The \texttt{InsightEngine} (\texttt{core/insight\_engine.py}) implements a continuous surprise-detection system over the reasoning stream. It fires an \texttt{InsightEvent} when a traversal path scores signi\-ficantly above the rolling baseline for its community pair.

\textbf{Surprise signal:}

\begin{equation}
\text{surprise}(P) = \text{score}(P) - \mu_{\text{baseline}}(C_u, C_v)
\end{equation}

where $\mu_{\text{baseline}}$ is a per-community-pair rolling mean maintained in a fixed-size ring buffer (O(1) update). An \texttt{InsightEvent} fires when:

\begin{equation}
\text{surprise}(P) > \theta_{salience} = 0.35
\end{equation}

\textbf{Insight score:}

\begin{equation}
\text{insight\_score} = \min\!\left(1,\; \frac{\text{surprise} + \text{explanatory\_power}}{2}\right)
\end{equation}

\begin{equation}
\text{explanatory\_power}(C_u, C_v) = \frac{|C_u| \times |C_v|}{N(N-1)/2}
\end{equation}

The explanatory power term rewards insights that bridge large, struc\-turally disco\-nnected communities — high-value discoveries that span conceptual gaps.

\textbf{Three-tier archi\-tecture:}

\begin{center}
{\footnotesize
\begin{tabularx}{\columnwidth}{|>{\raggedright\arraybackslash}X|>{\raggedright\arraybackslash}X|>{\raggedright\arraybackslash}X|}
\hline
Tier & Mechanism & Cost \\ \hline
Hot path & O(1) ring-buffer surprise check per traversal step & ~0\% CPU \\ \hline
Warm path & Daemon thread, event queue drain + baseline update & ~0.01\% CPU \\ \hline
Cold path & Full community-boundary scan (scheduled) & ~50 ms/hr \\ \hline
\end{tabularx}
}
\end{center}

\textbf{Materi\-alization:} When an \texttt{InsightEvent} fires, the engine mater\-ializes an \texttt{INSIGHT\_LINK} edge on the insight path:

\begin{equation}
\text{INSIGHT\_LINK}: \text{confidence} = 0.85,\; w = 2.0,\; \text{provenance} = \text{"insight"}
\end{equation}

Because \texttt{confidence = 0.85 \textbackslash gg \textbackslash theta\_{prune} = 0.2}, insight links \textbf{resist REM pruning} by a wide margin.

\textbf{Hebbian reward propagation:} Every edge on the insight path receives a confidence boost:

\begin{equation}
\Delta\text{confidence}(e) = \delta_{hebbian} \times \text{insight\_score}, \quad \text{capped at } 1.0
\end{equation}

\textbf{REM + Insight feedback loop:}

\begin{equation}
\underbrace{\text{REM proposes}}_{\text{confidence} = 0.3} \xrightarrow{\text{traversal validates}} \underbrace{\text{Insight cements}}_{\text{confidence} = 0.85} \xrightarrow{\text{resists}} \underbrace{\text{next REM prune}}_{\theta_{prune} = 0.2}
\end{equation}

This loop implements a structural analog of the dopamine prediction-error signal + AMPA receptor upreg\-ulation cycle in biological memory conso\-lidation.

\textbf{API endpoints:}

\begin{center}
{\footnotesize
\begin{tabularx}{\columnwidth}{|>{\raggedright\arraybackslash}X|>{\raggedright\arraybackslash}X|>{\raggedright\arraybackslash}X|}
\hline
Endpoint & Method & Description \\ \hline
\texttt{/insight/status} & GET & Engine state, event counts, baseline stats \\ \hline
\texttt{/insight/events} & GET & Recent \texttt{InsightEvent} list \\ \hline
\texttt{/insight/scan} & POST & Trigger immediate cold-path community-boundary scan \\ \hline
\end{tabularx}
}
\end{center}

Biological analog: \textbf{dopamine prediction-error signal} (surprise fires when actual reward exceeds expected) combined with \textbf{AMPA receptor upreg\-ulation} (Hebbian boost to the edges that carried the surprising path).

---

\subsection{12.10 Insight Validation and Metaco\-gnition (Phase 16)}

Phase 16 adds two components that operate above the InsightEngine: a structural verifier for individual discoveries and a second-order pattern detector across the discovery history.

\subsubsection{12.10.1 InsightValidator — Bilateral Verifi\-cation and Triang\-ulation}

A raw \texttt{InsightEvent} represents a \textit{surprising} connection. Surprise is a necessary but insuf\-ficient condition for a discovery to be trusted. \texttt{InsightValidator} (\texttt{core/insight\_validator.py}) applies two independent structural tests before promoting an insight's confidence.

\textbf{Check 1 — Bilateral reverse traversal:}

For an insight connecting source $s$ to target $t$, the bilateral check asks whether the \textit{original graph structure} (excluding the INSIGHT\_LINK edge itself and the direct $s \leftrightarrow t$ edge) can indep\-endently reach $s$ from $t$ within max\_hop steps:

\begin{equation}
\text{bilateral}(s, t) = \exists\; \text{path}_{G \setminus \{\text{INSIGHT\_LINK}, (s,t)\}}(t \to s,\; \text{length} \leq h_{\max})
\end{equation}

The exclusion of both the INSIGHT\_LINK edge and the direct pair prevents circular reasoning: the discovered connection cannot serve as its own evidence. If neither the insight link nor the original cross-edge can be used, a return path found in the remaining structure constitutes genuine independent confi\-rmation.

\textbf{Check 2 — Multi-path corro\-boration (trian\-gulation):}

Up to $N_{seeds}$ (default 5) other nodes in the same source community as $s$ are tested as independent starting points. The source node $s$ itself is removed from the graph to prevent corro\-borating paths from simply routing through $s$'s existing connection:

\begin{equation}
\text{corroboration}(t) = \left|\left\{ u \in \mathcal{C}_s \setminus \{s, t\} \;\middle|\; \exists\; \text{path}_{G_{\text{und}} \setminus \{s\}}(u \to t,\; \text{length} \leq h_{\max}) \right\}\right|
\end{equation}

If multiple community-peers indep\-endently reach $t$, the structural robustness of the connection is confirmed — it is not an artifact of one unusual path from $s$.

\textbf{Validation outcomes:}

\begin{center}
{\footnotesize
\begin{tabularx}{\columnwidth}{|>{\raggedright\arraybackslash}X|>{\raggedright\arraybackslash}X|>{\raggedright\arraybackslash}X|}
\hline
Status & Condition & Confidence after validation \\ \hline
\texttt{corroborated} & bilateral ∧ corro\-boration $\ge$ 2 & 0.95 \\ \hline
\texttt{bilateral} & bilateral ∧ corro\-boration \textless  2 & 0.92 \\ \hline
\texttt{unilateral} & ¬bilateral ∧ corro\-boration $\ge$ 1 & unchanged (0.85) \\ \hline
\texttt{isolated} & ¬bilateral ∧ corro\-boration = 0 & unchanged (0.85), flagged \\ \hline
\end{tabularx}
}
\end{center}

Confidence promotion updates the INSIGHT\_LINK edge in the graph so future traversal prefe\-rentially uses validated paths.

\textbf{API endpoints:}

\begin{center}
{\footnotesize
\begin{tabularx}{\columnwidth}{|>{\raggedright\arraybackslash}X|>{\raggedright\arraybackslash}X|>{\raggedright\arraybackslash}X|}
\hline
Endpoint & Method & Description \\ \hline
\texttt{/insight/validate/all} & POST & Validate all pending InsightEvents \\ \hline
\texttt{/insight/validate/{event\_id}} & POST & Validate a single InsightEvent by ID \\ \hline
\end{tabularx}
}
\end{center}

\subsubsection{12.10.2 MetaInsightEngine — Second-Order Reasoning}

The \texttt{MetaInsightEngine} (\texttt{core/meta\_insight\_engine.py}) watches the stream of \texttt{InsightEvent} objects and maintains an \textbf{InsightGraph}: a directed NetworkX graph where nodes are InsightEvent IDs and edges are structural relat\-ionships between insights.

\textbf{Four connection types} are detected for every pair of insights $(A, B)$:

\begin{center}
{\footnotesize
\begin{tabularx}{\columnwidth}{|>{\raggedright\arraybackslash}X|>{\raggedright\arraybackslash}X|>{\raggedright\arraybackslash}X|}
\hline
Type & Condition & Score \\ \hline
\texttt{chain} & $A.\text{target} = B.\text{source}$ or $B.\text{target} = A.\text{source}$ & $\frac{s_A + s_B}{2}$ \\ \hline
\texttt{shared\_entity} & $\{s_A, t_A, \text{bridge}_A\} \cap \{s_B, t_B, \text{bridge}_B\} \neq \emptyset$ & score $\times$ 0.8 \\ \hline
\texttt{community\_overlap} & $\text{leap}_A = \text{leap}_B \geq 1$ & score $\times$ 0.6 \\ \hline
\texttt{temporal\_cluster} & $|T_A - T_B| \leq \Delta t_{\text{window}}$ & score $\times$ 0.4 \\ \hline
\end{tabularx}
}
\end{center}

When an InsightGraph edge score exceeds \texttt{chain\_score\_threshold} (default 0.3), a \texttt{MetaInsightEvent} fires — the system has recognized a pattern in its own discovery history.

\textbf{Depth-2 higher-order detection:}

When a new insight $C$ arrives and finds prede\-cessors $B$ who themselves have prede\-cessors $A$, a depth-2 \texttt{MetaInsightEvent} fires:

\begin{equation}
A \xrightarrow{\text{meta}} B \xrightarrow{\text{meta}} C \implies \text{MetaInsightEvent}(A, C,\; \text{depth}=2,\; \text{chain}=[A, B, C])
\end{equation}

This corresponds to the recursive self-awareness expressed as "I notice that my discoveries about $B$ and $C$ are themselves connected to an earlier discovery about $A$." At depth 2, the system is reasoning about relat\-ionships between relat\-ionships — the structural equivalent of analogical reasoning.

The biological corre\-spondence is the formation of \textbf{episodic meta-memories}: the brain doesn't only remember individual experiences, it remembers that two experiences were related, and can later notice that this remembered relat\-ionship is itself related to a third. The InsightGraph is the structural substrate of this capacity in CEREBRUM.

\textbf{API endpoints:}

\begin{center}
{\footnotesize
\begin{tabularx}{\columnwidth}{|>{\raggedright\arraybackslash}X|>{\raggedright\arraybackslash}X|>{\raggedright\arraybackslash}X|}
\hline
Endpoint & Method & Description \\ \hline
\texttt{/meta-insight/status} & GET & Engine state, total meta-events, InsightGraph size \\ \hline
\texttt{/meta-insight/events} & GET & Recent \texttt{MetaInsightEvent} list \\ \hline
\texttt{/meta-insight/graph} & GET & Full InsightGraph as JSON (nodes + edges) \\ \hline
\end{tabularx}
}
\end{center}

\textbf{Biological analog:} Formation of \textbf{episodic memory clusters} in the hippocampus — the binding of related memories into schemata, and subsequent schema-to-schema assoc\-iations during memory conso\-lidation. Phase 16 closes the loop started in Phase 14 (REM offline conso\-lidation) and Phase 15 (online surprise detection): discoveries are now verified before being trusted, and the pattern of verified discoveries is itself analyzed for higher-order structure.

---

\subsection{13. Conclusion}

We have presented CEREBRUM: a framework that enables Knowledge Graphs to
reason using the structural principles of Transformer attention without
training data, without an LLM, and with full inter\-pretability.

The two core contr\-ibutions — Community-Structured Attention (CSA) and
Dual-Signal Community Fusion (DSCF) — work together to give a KG the dual
character of multi-head attention: local cohesion (from DSCF's LPA component)
combined with global structural signi\-ficance (from DSCF's modularity
component).

The resulting system produces reasoning paths, not \textbf{Black-Box} embeddings. Every
answer is traceable to a sequence of verified graph edges. This archi\-tectural shift 
moves AI from proba\-bilistic hidden-layer weights to a \textbf{Glass-Box} of deter\-ministic 
paths — a vital transition in the modern AI/ML landscape. Every reasoning step
names the community it traversed. This inter\-pretability property, combined with
the graph-grounded capability of graph-grounded inference, positions CEREBRUM
as a meaningful complement to — and in certain domains, replacement for —
LLM-based reasoning over structured knowledge.

The open questions identified in Section 10 define the ongoing research program.
The benchmarks in Section 9 define the empirical standard. The archi\-tecture in
Section 6 defines the core build. Sections 11–12 demonstrate that the
archi\-tecture scales to production, real-time streaming, experience-dependent
structural plasticity, autonomous causal discovery, offline memory conso\-lidation
(REM Cycle, Phase 14), surprise-driven insight discovery with Hebbian reward
propagation (InsightEngine, Phase 15), and bilateral verif\-ication + second-order
metac\-ognitive reasoning (InsightValidator + MetaInsightEngine, Phase 16) without
modif\-ication to the core CSA or DSCF algorithms.

The name CEREBRUM refers to the optical phenomenon where two viewpoints on
the same object yield depth perception that neither viewpoint alone provides.
LPA and modularity are two viewpoints on the same graph. Their combination
yields structural depth — attention heads with both short-range and long-range
character — that neither produces alone. This multi-signal consensus is inspired 
by \textbf{mid-level voting} systems in triplex-redundant aircraft navigation, where 
the median value is selected to correct navigation errors. CEREBRUM applies this 
principle to "right the navigation errors" (hallu\-cinations) of current language 
models by requiring structural consensus for every reasoning step.

That depth is what makes the KG reason.

---

\subsection{References}

\begin{enumerate}
\item Scarselli et al., "The Graph Neural Network Model," IEEE TNNLS, 2009.
\item Gilmer et al., "Neural Message Passing for Quantum Chemistry," ICML, 2017.
\item Velickovic et al., "Graph Attention Networks," ICLR, 2018.
\item Hamilton et al., "Inductive Repres\-entation Learning on Large Graphs," NeurIPS, 2017.
\item Bordes et al., "Translating Embeddings for Modeling Multi-relational Data (TransE)," NeurIPS, 2013.
\item Sun et al., "RotatE: Knowledge Graph Embedding by Relational Rotation in Complex Space," ICLR, 2019.
\item Xiong et al., "DeepPath: A Reinfo\-rcement Learning Method for Knowledge Graph Reasoning," EMNLP, 2017.
\item Das et al., "Go for a Walk and Arrive at the Answer (MINERVA)," ICLR, 2018.
\item Yao et al., "KG-GPT: A General Framework for Reasoning on Knowledge Graphs Using LLMs," 2023.
\item Chen et al., "KGPT: Knowledge-Grounded Pre-Training for Data-to-Text Generation," EMNLP, 2020.
\item Edge et al., "From Local to Global: A Graph RAG Approach to Query-Focused Summar\-ization," Microsoft Research, 2024.
\item Sarthi et al., "RAPTOR: Recursive Abstractive Processing for Tree-Organized Retrieval," ICLR, 2024.
\item Blondel et al., "Fast Unfolding of Communities in Large Networks (Louvain)," JSTAT, 2008.
\item Traag et al., "From Louvain to Leiden: promoting Well-Connected Communities," Scientific Reports, 2019.
\item Raghavan et al., "Near Linear Time Algorithm to Detect Community Structures in Large-Scale Networks (LPA)," Physical Review E, 2007.
\item Galarraga et al., "AMIE: Association Rule Mining under Incomplete Evidence in Ontological Knowledge Bases," WWW, 2013.

\end{enumerate}
---

\subsection{Appendix A: TSC Algorithm — Full Pseudocode}

\begin{lstlisting}
FUNCTION tsc_communities(G, resolution=1.0, max_iter=100,
                          temp_start=1.0, cooling=0.92,
                          centrality_weights=None):

  m = |E(G)|
  IF m == 0: RETURN [{v} for v in V(G)]

  nodes  = list(V(G))
  degree = {v: deg(v) for v in nodes}

  assignment = {v: i for i, v in enumerate(nodes)}   // singleton init
  com_k = {i: degree[v] for i, v in enumerate(nodes)} // degree-sum cache

  temperature = temp_start

  FOR iteration = 1 to max_iter:
    changed = False
    SHUFFLE nodes

    FOR EACH v in nodes:
      neighbors = N(v)
      IF |neighbors| == 0: CONTINUE

      cur = assignment[v]; kv = degree[v]

      // LPA signal
      vote = COUNT(assignment[nb] for nb in neighbors)
      lpa_cid = argmax(vote); lpa_conf = vote[lpa_cid] / |neighbors|

      // Modularity signal
      candidates = {assignment[nb] for nb in neighbors} - {cur}
      best_cid = cur; best_dq = 0
      FOR cid IN candidates:
        k_vc = |{nb in neighbors : assignment[nb] == cid}|
        dq = k_vc/m - resolution × kv × com_k[cid] / (2m²)
        IF dq > best_dq: best_dq = dq; best_cid = cid
      mod_conf = min(best_dq × m, 1.0)

      // Decision
      IF lpa_cid == best_cid != cur:
        new = lpa_cid  // consensus anchor
      ELIF best_cid == cur AND lpa_cid == cur:
        CONTINUE       // both say stay
      ELIF best_cid == cur:
        IF random() >= lpa_conf × temperature: CONTINUE
        new = lpa_cid
      ELIF lpa_cid == cur:
        IF random() >= mod_conf × (1 + (1−temperature)): CONTINUE
        new = best_cid
      ELSE:
        lpa_w = lpa_conf × temperature
        mod_w = mod_conf × (2 − temperature)
        new = weighted_choice({lpa_cid: lpa_w, best_cid: mod_w})

      com_k[cur] = max(com_k[cur] − kv, 0)
      com_k[new] = com_k[new] + kv
      assignment[v] = new; changed = True

    temperature = max(temperature × cooling, 0.01)
    IF NOT changed: BREAK

  // Connectivity post-pass (Leiden-style)
  // First component keeps original ID; additional components get new IDs.
  // (Preserves ID stability for majority partition; important for
  //  community_score lookup caching.)
  RETURN [component for community in assignment.values()
          for component in connected_components(G.subgraph(community))]
\end{lstlisting}

---

\subsection{Appendix B: CSA Weight Formula — Parameter Sensitivity}

The default parameter values ($\alpha$=0.4, $\beta$=0.4, $\gamma$=0.1, $\delta$=0.05, $\epsilon$=0.05) were
chosen based on the following intuitions:

\begin{itemize}
\item Embedding similarity ($\alpha$) and community membership ($\beta$) are given equal weight
\end{itemize}
  because both capture compl\-ementary aspects of relevance: similarity captures
  semantic proximity while community captures structural proximity.
\begin{itemize}
\item Edge type ($\gamma$) is given lower weight because it is most useful in domain-
\end{itemize}
  specific settings with rich edge type vocab\-ularies.
\begin{itemize}
\item Distance penalty ($\delta$) is kept small to allow multi-hop paths without excessive
\end{itemize}
  pruning.
\begin{itemize}
\item Hop decay ($\epsilon$) is minimal to allow deep traversal when needed.

\end{itemize}
In practice, $\alpha$ + $\beta$ dominate the attention weights for most graphs.

---

\subsection{Appendix C: Relati\-onship to Existing KG Embedding Methods}

TransE [Bordes et al., 2013] represents relations as trans\-lations in embedding
space: emb(h) + emb(r) $\approx$ emb(t) for (h, r, t) triples. CEREBRUM can use
TransE embeddings directly for the similarity term in CSA without modif\-ication.

RotatE [Sun et al., 2019] represents relations as rotations in complex space.
More expressive for symmetric, antis\-ymmetric, and compo\-sitional relations.
Also directly usable in CEREBRUM.

Neither TransE nor RotatE produces multi-hop reasoning paths on their own.
CEREBRUM uses their embeddings as the semantic grounding layer while the
traversal logic and community structure provide the reasoning.

---

\subsection{Appendix D: Prototype Code (Deprecated)}

\textgreater  \textbf{Note}: The prototype code previously contained in this appendix has been
\textgreater  superseded by the production imple\-mentation in \texttt{core/community\_engine.py}.
\textgreater  Refer to \textbf{Appendix A} for the current TSC algorithm pseudocode.

---

\subsection{Appendix E: Project Bootstrap Guide}

\subsubsection{E.1 Phase 0 Completion Status}

Phase 0 is \textbf{complete}. The following work is done:

\begin{itemize}
\item [x] White paper written (Sections 1-11, Appendices A-C, this document)
\item [x] DSCF algorithm designed and formalized
\item [x] DSCF prototype implemented in Python (Appendix D.1)
\item [x] Leiden/LPA wrappers implemented (Appendix D.2)
\item [x] DSCF validated on Home Assistant's Neo4j graph (stable communities, correct behavior)
\item [x] CSA attention formula designed and documented
\item [x] Transformer-to-KG structural equivalence table complete
\item [x] Full repo structure designed
\item [x] Hypotheses H1, H2, H3 stated and evaluation protocol defined
\item [x] Benchmark datasets identified (WebQSP, MetaQA-2hop/3hop, FB15k-237)

\end{itemize}
\subsubsection{E.2 Creating the Standalone Repo}

\begin{lstlisting}
# From E:\Development\ (or wherever you keep projects)
mkdir parallax
cd parallax
git init
git commit --allow-empty -m "chore: initial repo"

# Copy this file as the living research document
cp ../Home Assistant/PARALLAX.md PAPER.md
git add PAPER.md
git commit -m "docs: add white paper v0.1 — Phase 0 complete"

# Create directory structure
mkdir -p core reasoning adapters llm_bridge api cli tests/fixtures benchmarks examples
touch core/__init__.py reasoning/__init__.py adapters/__init__.py

# Start Phase 1: extract community_engine.py from Appendix D
# Copy D.1, D.2, D.3 code blocks into:
#   core/community_engine.py   (dscf_communities, leiden_communities, lpa_communities, hybrid_communities)
#   core/structural_encoder.py (compute_structural_features)
\end{lstlisting}

\subsubsection{E.3 Depend\-encies}

\textbf{Core (required):}
\begin{lstlisting}
networkx>=3.0.0
numpy>=1.24.0
igraph>=0.10.0
leidenalg>=0.10.0
scipy>=1.10.0
\end{lstlisting}

\textbf{Embeddings (choose one):}
\begin{lstlisting}
sentence-transformers>=2.2.0    # Option B: label-based, zero training (recommended default)
pykeen>=1.10.0                  # Option A: TransE/RotatE structural embeddings
\end{lstlisting}

\textbf{API (optional):}
\begin{lstlisting}
fastapi>=0.100.0
uvicorn[standard]>=0.23.0
pydantic>=2.0.0
\end{lstlisting}

\textbf{Graph DB adapters (optional):}
\begin{lstlisting}
neo4j>=5.8.0        # Neo4j adapter
SPARQLWrapper>=2.0  # RDF/Wikidata adapter
\end{lstlisting}

\textbf{pyproject.toml} (starter):
\begin{lstlisting}
[build-system]
requires = ["setuptools>=68", "wheel"]
build-backend = "setuptools.backends.legacy:build"

[project]
name = "parallax-kg"
version = "0.1.0"
description = "Community-Structured Graph Attention for Knowledge Graph Reasoning"
authors = [{name = "Bryan Alexander Buchorn", email = "bryan.alexander@buchorn.com"}]
license = {text = "Proprietary"}
requires-python = ">=3.10"
dependencies = [
    "networkx>=3.0.0",
    "numpy>=1.24.0",
    "igraph>=0.10.0",
    "leidenalg>=0.10.0",
    "scipy>=1.10.0",
]

[project.optional-dependencies]
embeddings = ["sentence-transformers>=2.2.0"]
kge = ["pykeen>=1.10.0"]
api = ["fastapi>=0.100.0", "uvicorn[standard]>=0.23.0", "pydantic>=2.0.0"]
neo4j = ["neo4j>=5.8.0"]
all = ["parallax-kg[embeddings,api,neo4j]"]
\end{lstlisting}

\subsubsection{E.4 Phase 1 Task Checklist}

\begin{lstlisting}
[ ] core/community_engine.py
    [ ] Copy dscf_communities() from Appendix D.1
    [ ] Copy leiden_communities(), lpa_communities(), hybrid_communities() from D.2
    [ ] Add modularity_score(G, communities) utility function

[ ] core/structural_encoder.py
    [ ] Copy compute_structural_features() from Appendix D.3
    [ ] Add encode_structural_features(features, embedding_dim) -> np.ndarray
        (concatenate + project PageRank/betweenness/degree into d-dim vector)

[ ] core/graph_adapter.py
    [ ] Define Entity and Edge dataclasses
    [ ] Define GraphAdapter ABC with 4 abstract methods (Section 6.3)
    [ ] Implement NetworkXAdapter as the default in-memory backend

[ ] core/embedding_engine.py
    [ ] Define EmbeddingEngine ABC
    [ ] Implement SentenceEngine (sentence-transformers, label-based)
    [ ] Implement RandomEngine (np.random, for unit tests)

[ ] core/attention_engine.py
    [ ] Implement CSAEngine with compute_weight(u, v, k, ...) -> float
    [ ] Implement community_score(u, v, communities) -> float
    [ ] Precompute community_distance matrix via BFS on community graph

[ ] tests/fixtures/toy_graph.csv
    [ ] ~200 nodes, ~400 edges, 3-4 natural communities
    [ ] Suitable for unit tests at all hop depths 1-4

[ ] tests/test_dscf.py
    [ ] test_singleton_init
    [ ] test_two_cliques_separate (obvious community structure)
    [ ] test_disconnected_components_split (post-pass check)
    [ ] test_determinism_with_seed
    [ ] test_convergence_within_max_iter

[ ] tests/test_csa.py
    [ ] test_same_community_weight_is_highest
    [ ] test_cross_community_decay_with_distance
    [ ] test_parameter_defaults_sum_to_one
\end{lstlisting}

\subsubsection{E.5 Relati\-onship to Home Assistant}

CEREBRUM is archi\-tecturally independent from Home Assistant. The only code shared is the
DSCF prototype (now extracted above). When CEREBRUM matures:

\begin{itemize}
\item Home Assistant's \texttt{knowledge\_service} can optionally import \texttt{parallax} as a library
\end{itemize}
  and replace its current community detection with \texttt{from parallax.core.community\_engine import dscf\_communities}
\begin{itemize}
\item The \texttt{neo4j\_adapter} in CEREBRUM will mirror patterns already in Home Assistant's \texttt{knowledge\_service}
\item Home Assistant's holographic memory WebSocket pattern can serve as reference for
\end{itemize}
  CEREBRUM's optional real-time community broadcast feature

No Home Assistant code other than Appendix D functions should be copied into CEREBRUM.

---

\textbf{Copyright © 2026 Bryan Alexander Buchorn. All Rights Reserved.}
This document and the software it describes are protected by inter\-national copyright laws. Unauth\-orized commercial repro\-duction, distr\-ibution, or use without express written permission is strictly prohibited.





\bibliographystyle{IEEEtran}
\bibliography{../../docs/latex/references}

\end{document}
